\documentclass[11pt]{article}
\usepackage[a4paper,top=2.4cm,bottom=2.5cm,left=2.0cm,right=2.0cm]{geometry}
\usepackage{amsmath,amssymb,amsfonts}
\usepackage{fontspec}
\setmainfont{TeX Gyre Termes}
\usepackage{booktabs}
\usepackage{graphicx}
\usepackage{pdflscape}
\graphicspath{{./}}
\usepackage[section]{placeins}
\usepackage{float}
\renewcommand{\topfraction}{0.92}
\renewcommand{\bottomfraction}{0.8}
\renewcommand{\textfraction}{0.07}
\renewcommand{\floatpagefraction}{0.75}
\usepackage{xcolor}
\usepackage{tikz}
\usetikzlibrary{arrows.meta,positioning}
\usepackage{caption}
\captionsetup{font=small,labelfont=bf}
\usepackage{titlesec}
\titleformat{\section}{\normalfont\large\bfseries}{\thesection}{0.6em}{}
\titleformat{\subsection}{\normalfont\normalsize\bfseries}{\thesubsection}{0.5em}{}
\titleformat{\subsubsection}{\normalfont\normalsize\bfseries}{\thesubsubsection}{0.5em}{}
\titlespacing*{\section}{0pt}{1.3ex plus .2ex minus .2ex}{0.6ex}
\titlespacing*{\subsection}{0pt}{1.0ex plus .2ex minus .2ex}{0.5ex}
\usepackage{fancyhdr}
\pagestyle{fancy}
\fancyhf{}
\fancyhead[L]{\small\itshape Experimental Astronomy}
\fancyhead[R]{\small\thepage}
\renewcommand{\headrulewidth}{0.4pt}
\usepackage{cite}
\usepackage{hyperref}
\hypersetup{hidelinks}
\setlength{\emergencystretch}{3em}
\newcommand{\keV}{\ensuremath{\,\mathrm{keV}}}
\newcommand{\MeV}{\ensuremath{\,\mathrm{MeV}}}
\newcommand{\cms}{\ensuremath{\,\mathrm{cm^{-2}\,s^{-1}}}}
\newcommand{\phcms}{\ensuremath{\,\mathrm{ph\,cm^{-2}\,s^{-1}}}}
\newcommand{\cps}{\ensuremath{\,\mathrm{s^{-1}}}}
\newcommand{\wii}{W_{511}}
\newcommand{\aeff}{A_{\mathrm{eff}}}
\newcommand{\ftoaref}{F_{\mathrm{TOA,ref}}}
\newcommand{\fpayloadref}{F_{\mathrm{payload,15,ref}}}
\newcommand{\fthree}{F_{3\sigma}}
\raggedbottom

\begin{document}
\thispagestyle{fancy}
\begin{flushleft}
{\LARGE\bfseries Detector-coupled Monte Carlo estimate of background and unresolved-line sensitivity for a balloon-borne 511 keV Laue-lens TES telescope\par}
\vspace{1.1em}
{\large Authors omitted for review\par}
\vspace{0.3em}
{\normalsize\itshape Affiliations omitted for review\par}
\end{flushleft}
\vspace{0.7em}
\noindent\rule{\textwidth}{0.4pt}
\vspace{0.5em}

\noindent\textbf{Abstract}\hspace{0.7em}
Detecting compact 511 keV emission with a balloon-borne Laue-lens telescope
requires prompt and activation backgrounds to be evaluated with the focused
signal under a common detector-response and event-selection framework. We use
an end-to-end Geant4/MEGAlib chain to transport a focused celestial signal,
seven prompt particle families, the broadband atmospheric $\gamma$ field, and
production-position-sampled activation decays through two complete
detector--cryostat mass models. All streams are subjected to the same
$510.58$--$511.42\keV$ science window, common-time event building, detector
response, active anticoincidence, and topology/aperture selection within each
design. Mass model A gives a day-15 direct, post-selection background rate of
$5.3057\times10^{-2}\cps$. In two independent weighted-rate projections of its
selected delayed component, the dilution-refrigerator/mixing-chamber and staged
cold structures account for 32.18\%, whereas copper accounts for 70.87\% by
material.
This provenance motivates evaluation of mass model B as a complete alternative
configuration. Model B retains a selected effective area of
$15.08544\pm0.04503\,\mathrm{cm^2}$, or 99.75\% of the model-A value, while
reducing the day-15 direct, post-selection background to $1.0747\times10^{-2}\cps$.
For the IBIS-based $3\sigma$-equivalent benchmark flux,
$\ftoaref=2.4\times10^{-4}\phcms$ at the top of
the atmosphere, equivalent to $\fpayloadref\simeq1.565\times10^{-4}\phcms$ at
the payload on day 15, an 81-node, 20 d fold gives model B 4028.15 signal and
18568.89 background counts, with Gaussian and Asimov significances of 29.561
and 28.578. The corresponding $3\sigma$ top-of-atmosphere thresholds are
$(2.436\pm0.223)\times10^{-5}$ and
$(2.445\pm0.223)\times10^{-5}\phcms$. Under the stated source, response,
timing, and selection conditions, model B reduces the 20 d background to
21.19\% of the model-A value and improves the Gaussian flux threshold by a
factor of 2.227; attribution to any individual design change requires a
dedicated ablation study.
An exploratory two-stream reweighting also identifies quiet LEO,
lunar-regolith, and Sun--Earth L2 source fields as distinct follow-up regimes;
the resulting screening indices are not environment-matched sensitivity
forecasts.

\vspace{0.7em}
\noindent\textbf{Keywords}\hspace{0.7em}Gamma-ray instrumentation, 511 keV annihilation line, transition-edge sensors, Laue lens, balloon-borne telescope, activation background

\vspace{0.3em}
\noindent\rule{\textwidth}{0.4pt}
\vspace{1.0em}

\section{Introduction}
\label{sec:introduction}

Galactic 511 keV annihilation radiation is a key observational probe of the
production, propagation, and annihilation environments of Galactic positrons.
Its spectrum contains a two-photon line near 511 keV and an ortho-positronium
continuum at lower energies. Positrons annihilate directly with electrons or
form para-positronium, producing two photons near 511 keV, whereas
ortho-positronium decays into a three-photon continuum. The line intensity,
profile, and positronium continuum therefore constrain the medium in which
positrons slow down and annihilate \cite{Prantzos2011,Jean2006}. Under an
approximate steady-state assumption, the total Galactic 511 keV flux measured
by INTEGRAL/SPI corresponds to an annihilation rate of a few
$\times10^{43}\,\mathrm{e^+\,s^{-1}}$ \cite{Prantzos2011,Kierans2020}. Proposed
positron sources include $\beta^+$-unstable nuclei synthesized in stellar
explosions, compact objects capable of producing electron--positron pairs, and
more speculative particle-physics scenarios in which annihilating or decaying
light dark matter injects low-energy positrons. Existing observations do not
uniquely determine the dominant contribution
\cite{Prantzos2011,Siegert2016AA,Siegert2016V404}.

Wide-field observations have established the diffuse morphology of the line.
INTEGRAL/SPI finds a bright central bulge and a weaker disk, and a recent
20-year analysis recovers central, broad-bulge, and disk components while the
evidence for additional correlated structures remains marginal
\cite{Knodlseder2005AllSky,Siegert2016AA,Yoneda2025}. The 2016 COSI balloon flight independently
detected the bulge and found evidence that the emission is more extended than a
single point source \cite{Kierans2020,Siegert2020COSI}. Spectroscopy also
constrains the annihilation medium and kinematics: the SPI spectrum can be
decomposed into a narrow component with about $1.3\keV$ FWHM and a broad
component with about $5.4\keV$ FWHM, and spatial variations in centroid energy
and width have been sought across the inner Galaxy
\cite{Jean2006,Siegert2019Kinematics}. The annihilation map, however, records
where positrons slow down and annihilate, not necessarily where they were
created. Propagation through a magnetized, multiphase interstellar medium means
that diffuse morphology and spectroscopy alone cannot recover the dominant
source class or the complete pre-annihilation transport history
\cite{Prantzos2011}.

Wide-field Compton and coded-mask instruments are well suited to the total
bulge and disk emission. A complementary pointed instrument is needed to test
whether the central region also contains an unresolved central enhancement,
compact sources, or measurable line-kinematic structure. A detection of a
point-like or narrowly concentrated feature would show that part of the
annihilation is spatially localized; a non-detection would constrain a compact
component without excluding source populations whose positrons escape and
annihilate diffusely. An INTEGRAL/IBIS compact-source search found no Galactic
511 keV point sources and set a $2\sigma$ upper limit of about
$1.6\times10^{-4}\phcms$ for a $\sim10$ Ms Galactic-centre exposure
\cite{DeCesare2011}. This limit sets both the comparison scale of earlier
compact-source searches and a sensitivity target for a new narrow-field
pointing instrument.

Laue optics offer one route to such pointed observations by using
transmission-mode Bragg diffraction to focus photons from a known direction
onto a small focal plane \cite{Barriere2009Laue,Barriere2011LauePrototype}.
Concentrating the flux collected over a large aperture onto a smaller sensitive
detector area can improve point-source sensitivity when detector-related
background dominates \cite{Weidenspointner2006MAX}. The focal plane must then
combine a narrow energy response with sufficient stopping efficiency at
511 keV. Transition-edge-sensor (TES) microcalorimeters use the steep
resistive transition of a superconductor to measure small temperature changes
from deposited energy and can provide sub-keV resolution at 511 keV
\cite{Shirazi2023,Zhang2022TESGamma}. An array covers the focal spot, while thick absorbers and a
multilayer layout raise the full-energy efficiency. Laue focusing compresses
the accepted phase space, and the TES response permits a narrow line window and
detailed line-profile measurements. The resulting instrument trades field of
view for a concentrated beam and high-resolution spectroscopy, complementing
rather than replacing wide-field bulge and disk measurements.

Energy resolution and focusing performance alone do not determine point-source
sensitivity \cite{Shirazi2023,Hu2026}. Prompt energy deposits from the radiation
environment and delayed emission from activated detector and cryostat materials
also set the background in the science window
\cite{Gallego2025Balloon,Tian2026DIXE}.

This work quantifies the high-altitude performance of a balloon-borne
Laue-lens TES-array telescope for a 511 keV point source at the Galactic center
and resolves the selected background by physical component, material, parent
nuclide, and production region.

\subsection{Science and instrument requirements}
\label{sec:requirements}

\begingroup
\clubpenalty=10000
\widowpenalty=10000

The simulated instrument is designed to detect either 511 keV point-source
emission or an unresolved central-excess component in the Galactic-center
region, neither of which has been detected by previous instruments.
INTEGRAL/SPI measurements show that the 511 keV sky is dominated by diffuse
bulge and disk emission \cite{Siegert2016AA,Yoneda2025}. De Cesare's dedicated
search for compact sources with INTEGRAL/IBIS found no Galactic 511 keV point
sources and reported a $2\sigma$ upper limit for a 10 Ms exposure toward the
Galactic center \cite{DeCesare2011}. To compare this limit with the $3\sigma$
sensitivity reported here, we linearly rescale the published limit under the
Gaussian approximation, obtaining an IBIS-based $3\sigma$-equivalent
top-of-atmosphere upper limit of $2.4\times10^{-4}\phcms$. At the altitude on
flight day 15, the remaining vertical atmospheric column depth is
$3.461\,\mathrm{g\,cm^{-2}}$. For the instrument axis, designed for a fixed
elevation of $45^\circ$, the plane-parallel approximation gives a slant column
depth of $4.895\,\mathrm{g\,cm^{-2}}$; using an effective dry-air mass
attenuation coefficient of $0.08736\,\mathrm{cm^2\,g^{-1}}$, consistent with NIST
XCOM data near 511 keV, gives a transmission of $T_{15,45}=0.6520$. We
therefore define a matched reference-flux pair rather than using one symbol for
two observing planes: $\ftoaref=2.4\times10^{-4}\phcms$ at the top of the
atmosphere and
$\fpayloadref=\ftoaref T_{15,45}\simeq1.565\times10^{-4}\phcms$ at the payload
on day 15. The mission calculation uses $\ftoaref$ and applies the
time-dependent atmospheric transmission once; $\fpayloadref$ is its day-15
equivalent. This pair provides the benchmark for the simulated signal counts
and significance, while minimum resolvable fluxes are reported at the top of
the atmosphere. Compared with the wide-field spectroscopic and
coded-mask searches conducted with INTEGRAL, the payload modeled here
introduces Laue focusing to improve the point-source response and uses a
high-energy-resolution, multilayer thick-absorber TES microcalorimeter array as
the focal-plane detector to improve its capability for 511 keV line
measurements \cite{Shirazi2023,Caputo2017}.

This objective imposes four requirements on the instrument and the simulation.
First, the telescope must focus the 511 keV photons collected by a large-aperture
lens onto a compact focal plane, thereby increasing the ratio of the collection
area of the focusing optics to the sensitive area of the detector; this is a
narrow-field pointed configuration, not a means of increasing sky coverage.
Second, the focal-plane detector must provide a sub-keV energy response
sufficient to resolve fine line structure. Third, the detector and shielding
model must suppress atmospheric, activation, and side-entry backgrounds.
Fourth, the simulation chain must treat prompt background, delayed activation,
focused signal, vetoes, event topology, and mission-time variation consistently,
so that signal and background are evaluated under explicitly defined
detector-selection criteria.

The remainder of this paper is organized as follows. Section~\ref{sec:geometry}
presents the detector and cryostat mass models. Section~\ref{sec:framework}
describes the simulation framework, event selection, mission-time trajectory
fold, and the response-weighted environmental screening. Section~\ref{sec:results}
presents the balloon post-selection rates, background provenance, and
sensitivity. Section~\ref{sec:environment_screening} gives the explicitly
conditional transfer to orbital and lunar-surface source proxies.
Section~\ref{sec:discussion} discusses the interpretation and scope, and
Section~\ref{sec:conclusions} summarizes the conclusions.

\endgroup

\section{Detector and cryostat mass model}
\label{sec:geometry}

This study compares two complete detector--cryostat mass models: the reference
mass model A and mass model B, which was developed from the background analysis
of model A. Both contain the same six-layer TES core and a compact dilution
refrigerator (DR) with a mixing chamber, staged cold plates, thermal and magnetic
shields, copper supports, active scintillators, and cryogenic service structures
\cite{Shirazi2023,Gottardi2021TESReview}. They differ in the placement of the
focal plane and in the physical entrance, local tungsten structure, and active
shield arrangement. Figure~\ref{fig:mass_model} compares the two geometries with
the same camera and physical scale. Downstream calculations use the same Laue
focal-plane photon input, TES response, source definitions, science window,
time grouping, and Compton topology/aperture rule; active channels follow each
complete mass-model configuration.

The TES detector concept uses an absorber-coupled AlMn-film structure developed
in our laboratory. The transition temperature of AlMn can be tuned through its
composition and annealing process; the material is suitable for array fabrication
and has relatively low magnetic-field sensitivity
\cite{Xie2026AlMn,Vavagiakis2017AlMnMagnetic}. In the present conceptual design,
the TES thermometer is an approximately $200\,\mathrm{nm}$-thick AlMn film with
a nominal Mn atomic concentration of $2000\,\mathrm{ppm}$, whereas the $\gamma$-ray absorber is a
millimetre-scale metal body. The absorber consists of
$1.5\times1.5\times3.0\,\mathrm{mm^3}$ Ta pixels; the high atomic number of Ta
increases the interaction efficiency for 511 keV photons. At 511 keV, a single
$3.0\,\mathrm{mm}$ Ta layer has an estimated interaction efficiency of about
$49\%$, and the corresponding six-layer efficiency is close to unity. In the
current geometry, adjacent Ta array layers have a center-to-center spacing of
$12\,\mathrm{mm}$; this spacing was chosen with a subsequent veto based on
Compton-kinematic track reconstruction in mind. Ta was also selected for its low
specific heat as a superconductor at cryogenic temperatures; the value used here
is $5.99\times10^{-7}\,\mathrm{J\,kg^{-1}\,K^{-1}}$. Ta also has a relatively
high photoelectric-to-Compton interaction ratio. From
the NIST XCOM cross sections near 511 keV, this ratio is approximately 0.731
\cite{NISTXCOM}. In the detector-response model, we use our laboratory experience
with a previously developed X-ray TES \cite{Xie2026AlMn} to adopt
$\Delta E_{\mathrm{FWHM}}\simeq420\,\mathrm{eV}$ at 511 keV as the estimated
energy resolution and detector-response post-processing input.

Because the TES film and Ta absorber differ greatly in dimensions and mass, and
because the TES film itself contributes negligibly to interactions at 511 keV,
the transport geometry uses the Ta absorber as the principal sensitive volume of
each TES pixel. The explicitly modelled sensitive volume comprises six layers,
each containing 376 Ta absorber pixels, for a total of 2256 pixels. Each pixel is
$1.5\,\mathrm{mm}$ long and wide and $3\,\mathrm{mm}$ thick. A Si substrate proxy
is placed below each Ta pixel; because detailed thermal conduction is not modelled,
Si represents the $\mathrm{SiO_2}$--$\mathrm{SiN_x}$--Si substrate. The TES film,
wiring, and readout details are not constructed layer by layer as explicit
interaction volumes, but enter the model through the assumed energy response and
service-mass proxies. Defined independently of the response FWHM, the retained science-window
analysis selection uses the 511 keV line window
$W_{511}=511\,\mathrm{keV}\pm420\,\mathrm{eV}$, or
$510.58$--$511.42\,\mathrm{keV}$.

The TES array in mass model A is mounted on copper supports immediately below the
staged cold plates. Along the optical path are successive aluminium thermal-shield
windows, a beryllium vacuum window, and a tungsten multihole collimator aligned
with the TES pixels. Boolean apertures through the surrounding shells define the
side-entry beam path, and the complete instrument frame is rotated by
$45^\circ$ so that the local aperture points $45^\circ$ above the horizon.

Mass model B uses the same six TES layers and Laue focal-plane photon input but
translates the focal plane laterally into a cylindrical aluminium chimney. The
chimney has inner and outer radii of 10.75 and 11.05 cm, a 3 mm wall, and a
local axis at $z'=-2.80$ cm. Its physical entrance has a radius of 2.70 cm, and
an open four-sided tungsten frame replaces the multihole collimator used in
model A. Active BGO covers the chimney side, the front optical annulus, and the
rear cold-port annulus.

\begin{figure}[t]
\centering
\includegraphics[width=0.98\linewidth]{figures/manuscript/fig01_mass_models_ab_en.pdf}
\caption{Matched detector--cryostat views of (a) mass model A and (b) mass
model B, rendered with the same camera, physical viewport, and scale. In model
A the six-layer TES array lies below the staged cold structures and behind a
multihole W collimator. Model B translates the same TES core into a lateral Al
chimney and uses an open four-sided W frame with active BGO on the chimney side,
front optical annulus, and rear cold-port annulus. Red arrows show the nominal
focused 511 keV paths. Selected outer envelopes and the model-A plastic
scintillator are translucent or omitted for clarity; model B has no plastic
channel. The physical entrances shown here are distinct from the analysis disks
used by the topology veto.}
\label{fig:mass_model}
\end{figure}

The matched view makes the focal-plane displacement explicit. In mass model A,
a 4.796-mm-thick Bi upper half-cylinder is mounted inside the 2-mm Al cold-stage
cylinder, above the six Ta absorber planes and their copper thermal support. In
mass model B, the six planes instead lie along the lateral chimney axis, outside
the direct projection of the staged cold plates and inside the local BGO
assembly.

Mass models A and B retain the same six sensitive planes, pixel dimensions,
response model, Laue focal-plane input, and optical-axis convention. They are
not a single-variable geometry pair: the physical entrance and W structure also
differ, as Figure~\ref{fig:mass_model} shows. The results below therefore compare
two complete designs under a common source, response, timing, and selection
contract; the mass-model-B sensitivity is conditional on its present entrance
and open-frame acceptance.

\FloatBarrier

\section{Simulation framework}
\label{sec:framework}

Detector-coupled Monte Carlo transport through detailed mass models was used to calculate the response and background of the 511\,keV spectrometer \cite{Sturner2003}. Figure~\ref{fig:workflow} summarizes the end-to-end calculation in three physical branches: focused signal, prompt atmospheric background, and delayed activation background. The Laue optics and the detector--cryostat system were represented by separate mass models. The reference point source was first transported through the Laue-lens model; the positions and directions of photons reaching the focal plane were written to a MEGAlib EventList and used to initialize detector-side Monte Carlo transport. For each detector--cryostat mass model, this focused-signal input and all background sources were transported through the same complete model described in Section~\ref{sec:geometry}.

Prompt transport and activation-production transport were run in parallel from
the same eight-family EXPACS/PARMA broadband-field definition: seven charged-
particle or neutron families and the broadband atmospheric $\gamma$ field. In
the delayed-activation branch, BUILDUP transport recorded nuclide yields,
material volumes, and production positions. These records were used to
calculate the activity accumulated during the flight and to construct the
radionuclide inventory at each reference epoch; delayed decays were then
sampled at the recorded production positions and transported through the
detector--cryostat model.

The focused signal and the three background streams used in the rate budget---
seven prompt particle families, broadband atmospheric $\gamma$ rays, and
delayed activation---were processed with the same detector response,
active-veto criterion, and Compton-topology test within each mass model. The
selected background events were decomposed by incident particle, parent
nuclide, material, and production volume. This provenance first diagnoses mass
model A and then motivates evaluation of mass model B as a complete alternative
configuration. After both designs had been fixed, their selected signal and
component-resolved background rates were folded over the adopted trajectory to
obtain accumulated counts and unresolved-line sensitivity. The subsections
below define the source models, transport normalization, event selection, and
mission-time calculation.

\begin{landscape}
\centering
\vspace*{\fill}
\includegraphics[width=0.94\linewidth]{figures/manuscript/fig02_workflow_en_threebranch_20260830.pdf}
\captionof{figure}{End-to-end simulation and analysis workflow. The focused
signal, prompt atmospheric background, and delayed activation form three
physical branches. The background budget resolves the prompt branch
into seven particle families and broadband atmospheric $\gamma$ rays. All
streams undergo the common detector response, active veto, and Compton-topology
test within each mass model. Selected-event provenance motivates the complete
model-B configuration; the fixed designs are then evaluated on the common
mission timeline.}
\label{fig:workflow}
\vspace*{\fill}
\end{landscape}

\FloatBarrier

\subsection{Laue optics and detector-coupled signal transport}
\label{sec:optics}

The simulated Laue lens uses mosaic Ge(111) crystals \cite{Barriere2009Laue,Barriere2011LauePrototype}. The optical design has a 10 m focal length, a single-ring reference response at 511\keV{}, and an effective area of $20.1\,\mathrm{cm^2}$ at 511\keV, calculated as
\begin{equation}
  A_{\mathrm{eff}}(E)=A_{\mathrm{inc}}\frac{N_{\mathrm{fp}}(E)}{N_{\mathrm{inc}}(E)}
  \label{eq:laue_aeff}
\end{equation}
where $A_{\mathrm{inc}}$ is the Monte Carlo incident sampling area, $N_{\mathrm{inc}}(E)$ is the number of incident photons, and $N_{\mathrm{fp}}(E)$ is the number of photons satisfying the focusing/focal-plane interception condition.
The mosaic-crystal implementation provides the specified angular acceptance and is used consistently throughout the optical transport.

The Laue optical branch models both focusing and photon transport through the Ge crystals.

Diffraction is implemented as a discrete application-layer process in Geant4 rather than as an efficiency applied after transport \cite{Agostinelli2003,Allison2016}. At angular offset $|\Delta\theta|$, the tabulated diffraction probability $R$ is first normalized to the unabsorbed fraction, $R_{\mathrm{eff}}=R/(1-A)$, where $A$ is the photoelectric-absorption probability. It is then converted to an equivalent mean free path $\lambda_{\mathrm{diff}}$ through
$1-\exp(-\ell/\lambda_{\mathrm{diff}})=R_{\mathrm{eff}}$, where $\ell$ is the photon path length through the crystal step; hence $\lambda_{\mathrm{diff}}=-\ell/\ln(1-R_{\mathrm{eff}})$. Geant4 samples this process together with photoelectric absorption and Compton scattering, and the shortest sampled interaction length determines which process occurs first. A photon not selected for diffraction therefore continues under the standard electromagnetic physics and may be absorbed, scattered, or transmitted.

In this work, the diffraction probability is determined from a precomputed XOP/CRYSTAL rocking curve. The angular dependences of the Ge(111) reflectivity, transmissivity, and absorption are generated with the CRYSTAL/diff\_pat module of the XOP toolkit (called through xoppylib using crystallographic data from DABAX) \cite{SanchezDelRio2011XOP} and stored as a transport lookup table during preprocessing. At each transport step inside a crystal, the angular deviation from the Bragg condition is evaluated and the corresponding diffraction probability is obtained by interpolation. The 30 arcsec FWHM mosaic spread is implemented by applying a Gaussian angular perturbation to the ideal Bragg-plane normal. When diffraction occurs, the outgoing direction is determined from the perturbed plane normal, while the photon energy is conserved. This treatment follows existing Geant4 implementations of crystal Bragg reflection \cite{Guan2023BraggReflection,Reiazi2025G4BraggReflection} and is applied to the energy range considered here.

Here the reference signal is a monochromatic 511\keV{} point source. Its
normalization uses the matched pair $\ftoaref=2.4\times10^{-4}\phcms$ and
$\fpayloadref\simeq1.565\times10^{-4}\phcms$ at the day-15 payload plane
(Section~\ref{sec:requirements}), motivated by the dedicated INTEGRAL/IBIS
compact-source upper limit \cite{DeCesare2011} rather than by any single
observed object. It is constructed as a MEGAlib far-field source
\cite{Zoglauer2006}, with the detailed construction given in
Section~\ref{subsec:pointsource}. The signal is transported through the optics
branch and injected into the detector mass model as a parallel-beam source;
prompt cosmic-ray background and activation decays are transported directly in
the detector/cryostat mass model. The primary background budget is therefore
defined for the detector/cryostat branch, while the optics branch supplies the
transported focal-plane signal phase space.

\subsection{Source construction}
\label{sec:background}

The reference point-source EventList, the prompt atmospheric field, and the
delayed-decay source are constructed independently. For the background budget,
the prompt field is retained as seven non-$\gamma$ particle families and one
broadband atmospheric $\gamma$ component; activation production and delayed
decay remain separate transport stages. This organization preserves the
normalization and provenance needed to compare all selected background streams
under the common detector response and event selection.

\subsubsection{Prompt atmospheric source}
\label{subsec:prompt_source}

The prompt atmospheric field is derived from EXPACS/PARMA. PARMA3.0 is an
analytical model parameterized from PHITS atmospheric-shower calculations and
benchmarked against terrestrial cosmic-ray measurements, and EXPACS is its
public tabulation and software interface
\cite{Sato2015EXPACS,Sato2016PARMA4}. PHITS supplies the underlying
particle-transport physics: primary cosmic rays interact with air nuclei, the
secondary hadrons, leptons, photons, and neutrons are transported through the
atmospheric column, and the resulting fluences are parameterized by altitude,
geomagnetic cutoff, solar modulation, energy, and direction. In this work
EXPACS/PARMA provides the differential incident flux as a function of particle
species, kinetic energy, zenith angle, geographic position, altitude, cutoff
rigidity, and solar activity, and is used only to construct the external
radiation fields; detector response, vetoing, topology selection, and
mission-time counting are applied downstream. The reference source definitions
use latitude $34^\circ$, longitude $100^\circ$, altitude $38\,\mathrm{km}$,
vertical cutoff rigidity $R_c=11.6\,\mathrm{GV}$, and the 2025-08-31 solar
condition corresponding to $W=118.3$.

The prompt source includes photons, neutrons, electrons, positrons, protons,
alpha particles, and negative and positive muons. For each particle family the
EXPACS/PARMA differential spectrum is converted into a MEGAlib far-field area
source covering the full zenith range $\theta=0^\circ$--$180^\circ$ and full
azimuth range $\phi=0^\circ$--$360^\circ$. The zenith dependence is represented
by 20 differential equal-$\mu$ angular bins, where $\mu=\cos\theta$ and each
bin subtends $\Delta\Omega=0.6283185307\,\mathrm{sr}$: the first ten bins cover
down-going particles, the last ten up-going or albedo-like particles. Both
hemispheres are retained in the prompt transport and in the
activation-production transport, so upward albedo particles are not discarded
at source construction. Figure~\ref{fig:expacs_fullsphere_flux} shows the
resulting full-sphere flux field.

\begin{figure}[t]
\centering
\includegraphics[width=\linewidth]{figures/manuscript/fig03_corrected_kev_sources.pdf}
\caption{EXPACS/PARMA atmospheric particle field used for prompt transport and
activation production. (a) Fluxes of the eight particle families integrated over
energy and solid angle in the down-going and up-going (albedo) hemispheres.
(b) Integrated fluxes in the 20 equal-$\mu$ zenith bins ($\mu=\cos\theta$) for
each particle family. The environment is latitude $34^\circ$, longitude
$100^\circ$, altitude $38\,\mathrm{km}$, $R_c=11.6\,\mathrm{GV}$, and $W=118.3$.}
\label{fig:expacs_fullsphere_flux}
\end{figure}

The full-sphere integrated fluxes in the adopted atmospheric tables are 4.79966, 0.462239, 0.199002,
0.116981, 0.112301, 0.0114871, 0.00496893, and
$0.00557001\,\mathrm{cm^{-2}\,s^{-1}}$ for $\gamma$, $n$, $e^-$, $e^+$, $p$,
$\alpha$, $\mu^-$, and $\mu^+$, respectively. All families use 20 equal-$\mu$
angular bins and a 60 cm far-field start sphere. Minority families may be sampled
more heavily, but every family retains its own transported exposure,
\begin{equation}
  r_{\mathrm{event},j}=\left(\sum_i \mathcal{T}_{ij}\right)^{-1},
\end{equation}
where $\mathcal{T}_{ij}$ is the exposure of subsample $i$ for family $j$.
Consequently, over-sampling changes Monte Carlo variance but not the expected
physical rate.

The source energies passed to Cosima are expressed in keV. For non-alpha
families, kinetic energies tabulated in MeV are converted as
$E_{\mathrm{tot,keV}}=10^3E_{\mathrm{MeV}}$; the alpha-particle tables are given
in MeV per nucleon, so the total kinetic energy is
$E_{\alpha,\mathrm{keV}}=4\times10^3E_{\mathrm{MeV/nucleon}}$. The energy spectrum
in each angular bin is normalized to unit integral, and the bin-integrated flux
sets its absolute amplitude.

If $I_j(E,\mu)$ denotes the EXPACS/PARMA differential intensity for family $j$,
the flux carried by equal-$\mu$ bin $m$ is
\begin{equation}
 \Phi_{j,m}=2\pi\int_{\mu_m}^{\mu_{m+1}}\!\mathrm d\mu
             \int I_j(E,\mu)\,\mathrm dE,
 \qquad \Delta\mu=0.1,
 \label{eq:source_bin_flux}
\end{equation}
so every bin subtends $\Delta\Omega=2\pi\Delta\mu=0.62831853\,\mathrm{sr}$.
The first ten bins cover the down-going hemisphere and the second ten cover the
up-going or albedo hemisphere; azimuth is complete in both cases.

All eight particle families use a far-field area source of radius
$R=60\,\mathrm{cm}$. For each sampled direction, starting positions are drawn
uniformly from the projected disk perpendicular to that direction, and momenta
point into the mass model; the geometric rate factor is the projected area
$\pi R^2$. Together, the 20 bins cover $4\pi$ incident solid angle. The same
atmospheric-field definition is used separately for prompt transport and
activation-production calculations.

\subsubsection{Delayed activation source}
\label{subsec:delayed_source}

The delayed source is built from a separate activation-production transport. It
is not counted as immediate detector background; instead, it records the isotope,
material volume, incident family, and production position of each radioactive
product. Ground-state half-lives are checked against NUBASE2020
\cite{Kondev2021NUBASE}. For isotope $k$,
\begin{equation}
  A_k(t_{\mathrm{flight}})=P_k\left[1-\exp(-\lambda_k t_{\mathrm{flight}})\right],
  \qquad \lambda_k=\frac{\ln2}{t_{1/2,k}},
\end{equation}
where $P_k$ is the production rate and the reference state is
$t_{\mathrm{flight}}=15\,\mathrm{d}$.

The default MEGAlib/Cosima event record is insufficient to reconstruct the
spatial distribution of this delayed source. We therefore add a custom recording
hook to Geant4's \texttt{SteppingAction}, save the production coordinate when each
radioactive product is formed, and match the day-15 population to those records by
material volume, nuclide, and excitation state. The resulting source description
is indexed by incident family, material category, production volume, and
coordinate.

Retaining radionuclide production positions is essential because whether a
delayed particle escapes the surrounding material, triggers the geometry-specific
active anticoincidence, or passes the Compton-topology selection depends on its production
location. Activation-production records are matched to the day-15 nuclide
inventory by material volume, nuclide, excitation state, and coordinate, and decay
positions are sampled from these records with weights proportional to activity.
The sampled events are then normalized to the total activity of the corresponding
incident-particle family, preserving both the spatial distribution and the total
activity. The activation-source positions in the two mass models, together with
the parent nuclides that contribute to the TES 511 keV response window, are shown
in Fig.~\ref{fig:activation_source_positions}.

During the simulation, each incident-particle family transports $10^6$ delayed
decays, so each mass model generates $8\times10^6$ delayed triggers. After the
science-window, geometry-specific active-anticoincidence, and Compton-topology selections, 394 and 111
Monte Carlo events remain in mass models A and B, respectively. These events are
then normalized to physical rates using the day-15 activity and transported sample
size of each particle family and compared over the same physical observation time.
In mass model A, the day-15 activity separated
by incident-particle family is dominated by protons (724.2736 Bq), neutrons
(330.0210 Bq), and $\alpha$ particles (287.5411 Bq); the remaining contributions
are $\gamma$ (6.5923 Bq), $e^+$ (2.9537 Bq), $e^-$ (0.6679 Bq), $\mu^-$
(0.2479 Bq), and $\mu^+$ ($1.26\times10^{-9}$ Bq). The corresponding activities
in mass model B are 105.5262, 68.5316, 34.2836, 2.0320, 0.6754, 0.2663, 0.0684,
and 0.0040 Bq for $p$, $n$, $\alpha$, $\gamma$, $e^+$, $e^-$, $\mu^-$, and
$\mu^+$, respectively.

The final delayed events retain the parent-nuclide identity and production volume.
For activation-background events near 511 keV, the main contributors in mass
model A include $^{62}\mathrm{Cu}$ in the L2 copper support ring and the 50 mK
mixing-chamber plate. In mass model B, the largest contribution is
$^{62}\mathrm{Cu}$ in the L5 copper ring, followed by $^{62}\mathrm{Cu}$ in the
TES heat-sink ring; contributions from the W frame and Al shell are smaller.
These background origins and their relationship to the geometry are discussed
further below. Table~\ref{tab:background_source_model} summarizes the
source-construction hierarchy used by the two mass models.

\begin{table}[H]
\centering
\caption{Source layers used by both geometries. The table defines source
construction; detector response and selection are applied downstream.}
\label{tab:background_source_model}
\footnotesize
\begin{tabular}{p{0.21\linewidth}p{0.34\linewidth}p{0.37\linewidth}}
\toprule
Layer & Physical role & Implementation \\
\midrule
Atmospheric particle field & Prompt detector hits and activation production &
Eight families, 20 equal-$\mu$ bins per family, full sphere, 60 cm start radius;
total kinetic energy in keV. \\
Atmospheric $\gamma$ member & Prompt atmospheric photons &
The $\gamma$ member of the eight-family field, represented by 20 equal-$\mu$
bins spanning the full sphere. \\
Activation production & Family--nuclide--volume--position history &
The same atmospheric particle field, with NUBASE-checked ground-state decay data. \\
Production-position decay & Delayed particles emitted from recorded material coordinates &
$10^6$ decays per family; eight family-resolved decay samples per geometry,
merged into the delayed stream before common-time event building. \\
\bottomrule
\end{tabular}
\end{table}

\begin{figure}[p]
\centering
\includegraphics[width=0.96\linewidth]{figures/manuscript/fig_activation_positions_model_a_en.pdf}
\vspace{0.35em}

\includegraphics[width=0.96\linewidth]{figures/manuscript/fig_activation_positions_model_b_en.pdf}
\caption{Delayed activation-production positions in mass model A (top) and mass
model B (bottom). In each row, the left panel shows the full mass model and the
right panel enlarges the TES region. Blue points denote all transported delayed
activation-source positions; coloured points denote activation parent nuclides
that produce background in the TES $510.58$--$511.42$ keV response window.}
\label{fig:activation_source_positions}
\end{figure}

\FloatBarrier

\subsubsection{Reference point-source construction}
\label{subsec:pointsource}

The focused-signal stream uses an on-axis monochromatic 511\keV{} far-field
point source. Its reference normalization is $\ftoaref=2.4\times10^{-4}\phcms$
at the top of the atmosphere, equivalent to
$\fpayloadref\simeq1.565\times10^{-4}\phcms$ at the day-15 payload plane
(Section~\ref{sec:requirements}) \cite{DeCesare2011}.

In the optics calculation, photons arrive parallel to the optical axis and are
sampled uniformly over the incident aperture $A_{\mathrm{inc}}$ defined by
Equation~(\ref{eq:laue_aeff}). The focal-plane response is determined by the
30 arcsec mosaic spread, the XOP/CRYSTAL rocking curve, and the 10 m focusing
geometry. Three independent optics samples containing $1.5\times10^5$ incident
photons produce 37,256 diffracted focal-plane crossings, of which 37,194 fall
inside the $r=1.898\,\mathrm{cm}$ beryllium entrance window. The resulting spot
radii are $r_{90}=1.03\,\mathrm{cm}$ and $r_{99}=1.25\,\mathrm{cm}$. The position
and direction of each accepted crossing form a common focal-plane EventList that
is transported separately through each mass model.

For a top-of-atmosphere flux $F_{\mathrm{TOA}}$, the focused rate at the detector
entrance is
\begin{equation}
  R_{\mathrm{fp}}=F_{\mathrm{TOA}}\,A_{\mathrm{fp}}\,T_{\mathrm{atm},45}.
\end{equation}
The common EventList contains $N_{\mathrm{fp}}=37{,}194$ entries and represents
an optical effective area $A_{\mathrm{fp}}=20.08476\,\mathrm{cm^2}$. It
corresponds to the effective exposure
$T_{\mathrm{exp}}=N_{\mathrm{fp}}/[F_{\mathrm{TOA}}A_{\mathrm{fp}}
T_{\mathrm{atm},45}]$, so each entry has physical rate weight
$1/T_{\mathrm{exp}}$. For mass model $j$, the selected effective area is
$A_{\mathrm{eff},j}^{\mathrm{sel}}=A_{\mathrm{fp}}N_{\mathrm{sel},j}/N_{\mathrm{fp}}$.

At day 15, the atmospheric transmission is $T_{\mathrm{atm},45}=0.652034$. In
mass model A, 28,612 entries lie in the measured science window and 28,006 pass
the final Compton-topology selection, corresponding to
$A_{\mathrm{eff,A}}^{\mathrm{sel}}=15.12324\pm0.04492\,\mathrm{cm^2}$. The
corresponding counts in mass model B are 28,434 and 27,936, with
$A_{\mathrm{eff,B}}^{\mathrm{sel}}=15.08544\pm0.04503\,\mathrm{cm^2}$. After applying atmospheric
transmission and the day-15 accidental-coincidence survival, the signal rates at
$F_{\mathrm{TOA}}=\ftoaref$ are $2.29144\times10^{-3}\cps$ and
$2.35000\times10^{-3}\cps$ for mass models A and B, respectively. The selected
effective areas differ by $0.25\%$, so the side-chimney geometry preserves the
signal acceptance to this level.

These effective areas are on-axis reference values. Source offset and flight
attitude jitter modify the focused response through the finite angular acceptance
of the Laue lens and are treated as pointing systematics on the signal effective
area \cite{Barriere2009Laue,Weidenspointner2006MAX}. In the mission calculation,
the telescope optical axis is fixed at $45^\circ$ elevation and the focal-plane
EventList is injected along this axis through the lateral entrance window.

\subsection{Detector response, event selection, and veto definitions}
\label{sec:selection}

The common-time calculation maps the prompt atmospheric and delayed-activation
background catalogues onto one explicit time axis before the active-veto and
Compton selections are applied. The seven prompt particle families and the
broadband atmospheric $\gamma$ component retain their independent physical-rate
normalizations within the prompt stream. Focused-signal events are inserted as
conditional probes of accidental-coincidence loss rather than added to the
background budget.

\subsubsection{Time axis and rate normalization}
\label{subsec:time_axis}

In a stationary radiation environment, mutually independent particle events
with a constant mean arrival rate have Poisson-distributed counts in a fixed
time interval. Conditional on the event count, their arrival times are
uniformly distributed within that interval. This provides a natural temporal
representation for rate-normalized Monte Carlo catalogues: the catalogue
supplies the energy deposits, positions, and topologies, while the Poisson
process determines how often and when those events occur on the physical time
axis.

The radiation environment evolves along the balloon trajectory, but much more
slowly than the $1\,\mu\mathrm{s}$ coincidence time scale adopted in this study.
Within each trajectory node, the event streams can therefore be approximated as
independent Poisson processes with constant arrival rates. The focused signal,
prompt atmospheric background, and delayed-activation background are generated
by separate transport calculations and do not share a physical time axis,
whereas the geometry-specific active anticoincidence and Compton-topology selections depend on all
energy deposits within the same time window. Evaluating accidental coincidences
therefore requires these catalogues to be mapped onto a common time axis.

The three transported streams acquire per-record rate weights $r_{km}$ during
the source-normalization stage of Section~\ref{sec:background}, where
$k\in\{\mathrm{signal},\mathrm{prompt},\mathrm{delayed}\}$ labels the catalogue
and $m$ labels a transported record. The total rate of catalogue $k$ is
$R_k=\sum_m r_{km}$. For a reference interval of duration $T$, the number of
instances drawn from catalogue $k$ is
\begin{equation}
 N_k\sim\mathrm{Pois}(R_kT).
 \label{eq:poisson_superposition}
\end{equation}
Records are sampled with replacement with probabilities $r_{km}/R_k$ and are
assigned independent uniform times in $[0,T]$. The resulting sequence preserves
both the physical rate and the energy, direction, and detector-response
distribution of the event stream.

The three draws are merged and ordered by arrival time. Events separated by no
more than $\tau=1\,\mu\mathrm{s}$ are linked transitively into one coincidence
candidate group $c$. All energy deposits in the group are then subjected jointly
to the science-window, geometry-specific active-anticoincidence, and Compton-topology selections.
Figure~\ref{fig:poisson_time_axis_schematic} summarizes this construction.

\begin{figure}[H]
\centering
\includegraphics[width=\linewidth]{figures/manuscript/fig04a_poisson_time_axis_schematic_threebranch_20260830.pdf}
\caption{Schematic construction of the common-time coincidence axis. (a) Each
transported stream carries per-record rate weights and a total rate.
(b) Independent Poisson draws are merged and time-ordered on a common axis.
(c) Events separated by no more than $\tau=1\,\mu\mathrm{s}$ form one
coincidence candidate group. Colours identify catalogue membership; event
positions and densities are schematic. Prompt and delayed streams define the
background axis; focused signal is inserted as a conditional probe.}
\label{fig:poisson_time_axis_schematic}
\end{figure}

This calculation is performed at five reference times on days 0, 5, 10, 15,
and 20. At each reference time, the physical rates of the seven prompt particle
families, broadband atmospheric $\gamma$ rays, and delayed activation set the
Poisson draws used to construct a common background time axis for that
reference environment. Events separated by no more than $1\,\mu\mathrm{s}$ form
one coincidence group and jointly undergo the science-window, geometry-specific
active-anticoincidence, and Compton-topology selections. Focused-signal events are
also inserted individually into the corresponding background axis to quantify
signal loss from accidental coincidences.

Instrument performance is then evaluated along the long-duration balloon
trajectory, with the 20 d flight divided into 6 h intervals. Particle fluxes,
atmospheric transmission of the signal, and the activation inventory evolve with flight
time and position, thereby changing both the event density on the common time
axis and the accidental-coincidence probability. The background corrections
and signal-survival fractions obtained at the five reference times are applied
along the full trajectory. Section~\ref{sec:mission_time_fold} specifies the
trajectory model, temporal interpolation, and cumulative-count calculation.

\subsubsection{Detector response, science window, and weighted rates}
\label{subsec:detector_response}

Before the science-window and anticoincidence decisions, TES energy deposits in
each candidate group are combined by pixel. Multiple deposits in the same pixel
are summed to give its true deposited energy, and their energy-weighted position
represents the interaction location. The expected single-pixel energy resolution
is $\Delta E_{\mathrm{FWHM}}=420\,\mathrm{eV}$; TES measurement fluctuations are
approximated by a Gaussian response with this FWHM. Pixel hits with measured
energies below $0.300\keV$ are removed. Mass models A and B use the same response
and threshold settings.

The measured energies of all retained TES pixel hits are summed to give the total
measured TES energy of the event, which is used for the 511\,keV science-window
selection. The energy, position, layer, and multiplicity of each pixel hit remain
available for the subsequent Compton-topology selection. BGO deposits are
accumulated separately, are not included in the TES total, and enter the active
anticoincidence decision in the next subsection. The plastic-scintillator channels
in mass model A are treated in the same separate manner.

The primary analysis uses the common science-window condition
$510.58\leq E_{\mathrm{TES}}<511.42\keV$. This
high-containment reference window was fixed in advance from the expected energy
resolution. It is applied identically to the focused signal and every background
component and is held fixed between the two mass models, rather than being tuned
separately for each geometry.

The quoted flux thresholds apply to source lines whose intrinsic widths are
small relative to the detector energy resolution, such that the
response-convolved line remains adequately contained within this science
window. A substantially broadened source line requires convolution of the
source profile with the detector response and a new optimization of the
integration window.

The source catalogues have different physical fluxes, transport exposures, and
source-level reweighting. Each event therefore carries a physical-rate weight,
and the signal or background rate at a given selection stage is obtained by
summing the weights of events that pass that stage. Raw record counts quantify
Monte Carlo transport support rather than physical rate. Finite-transport
statistical uncertainties are propagated from the actual weight distribution,
while focused-signal acceptance is treated binomially using the common focal-plane
sample. Long-duration rate propagation along the mission trajectory is described
in Section~\ref{sec:mission_time_fold}; the response-convolved spectra through
the active-veto and topology/aperture stages are reported with the cut-flow results.

\subsubsection{Active anticoincidence gate}

Active anticoincidence rejects events that produce a 511\,keV candidate in the
TES while also depositing energy in the surrounding active shield. Passive
shielding attenuates incident radiation, whereas the BGO and plastic
scintillators also provide time-coincident information for event-by-event
discrimination. A focused 511\,keV photon reaches the TES mainly through the
optical channel. Atmospheric charged particles, secondaries from hadronic
interactions, and some broadband-$\gamma$ cascades are more likely to deposit
energy in both the TES and the surrounding active volumes. Coincident
active-shield energy can therefore identify background candidates in the
science window.

The active gate uses the common time axis and $1\,\mu\mathrm{s}$ coincidence
groups defined in Section~\ref{subsec:time_axis}. Within each candidate group,
TES deposits are first combined by pixel and then passed through the response
defined in Section~\ref{subsec:detector_response}; the science-window decision
uses the sum of the measured energies of all retained pixels. Simulated energy
deposits in all BGO volumes in the same group are summed as
$E_{\mathrm{BGO}}$. For mass model A, deposits in all plastic-scintillator
volumes are summed separately as $E_{\mathrm{pl}}$. These active-volume sums
are used only for anticoincidence and do not contribute to the TES event energy.

The active system of mass model A comprises a BGO side shell, bottom cap, and
top annulus together with plastic-scintillator side, bottom, and top sections.
A candidate group passes when
\[
 E_{\mathrm{BGO}}<50\keV
 \quad\text{and}\quad
 E_{\mathrm{pl}}<50\keV .
\]
Mass model B has no plastic scintillator. Its active system comprises the
chimney side shield, front optical annulus, and rear cold-port annulus, and its
passing condition is
\[
 E_{\mathrm{BGO}}<50\keV .
\]
Each threshold is applied to the corresponding active-material sum over the
entire coincidence group; a sum equal to or above $50\keV$ vetoes the group.
Within a given mass model, the same group-level criteria are applied to the
focused signal, seven-family prompt background, broadband atmospheric
$\gamma$ background, and delayed activation.

The groupwise sums include both correlated TES--shield deposits from one
incident particle and accidental overlap of otherwise independent hits whose
arrival times satisfy the $1\,\mu\mathrm{s}$ linking rule. A focused photon can
therefore be rejected when an unrelated active-shield hit enters the same
coincidence group. The associated loss is described by the accidental
signal-survival fraction defined in Section~\ref{subsec:time_axis} and included
in the mission-time fold. Candidates that pass the active gate proceed to the
Compton-topology and entrance-aperture consistency selection.

The analysis uses simulated energy deposits recorded by Cosima and applies
$50\keV$ as the common offline anticoincidence threshold.

\subsubsection{Compton topology and entrance-aperture consistency}

Active anticoincidence removes candidates with energy deposits above the adopted
threshold in the surrounding scintillators, but it cannot distinguish a focused photon from a
background event whose interactions remain within the TES array. Either can
leave a total energy near 511\,keV in the focal plane. A focused photon must,
however, arrive through the side-entry optical channel. For an event with two
or more separated TES hits, the division of energy among the hits and their
relative positions constrain the possible arrival directions. We therefore test
whether each event is compatible with the entrance aperture. The test uses the
kinematic cone construction of a Compton camera
\cite{Zoglauer2006}, but only as an aperture-consistency veto, not for imaging
or source localization.

The six-layer pixelated TES focal plane provides the information needed for
this test at pixel scale. For the retained hits defined in
Section~\ref{subsec:detector_response}, the response model supplies a measured
energy, while the event catalogue preserves the TES layer and pixel position.
Hits in different pixels or layers therefore provide both an energy division
and a finite spatial lever arm. The catalogue does not determine the interaction
order or the position within a pixel. A single-hit event consequently carries
no Compton-direction information and is retained as a calorimetric line
candidate; multi-hit events proceed to the consistency test below.

For an assumed interaction order, let a candidate of total energy $E_0$ deposit
$E_1$ at the first representative position $\mathbf r_1$ and then propagate
with $E_{\mathrm{sc}}=E_0-E_1$ to a second position $\mathbf r_2$. The measured
energies give the Compton scattering angle
\begin{equation}
  \cos\theta_C =
  1 - m_ec^2\left(\frac{1}{E_{\mathrm{sc}}}
  - \frac{1}{E_1+E_{\mathrm{sc}}}\right).
\end{equation}
The vector $\mathbf r_2-\mathbf r_1$ fixes the scattered-photon direction
$\hat{\mathbf u}_{12}$, but not the incident direction. Together,
$\hat{\mathbf u}_{12}$ and $\theta_C$ define a cone of allowed incident
directions with apex $\mathbf r_1$ and axis $-\hat{\mathbf u}_{12}$. The event
is aperture compatible if this cone intersects a circular analysis disk in the
side-entry plane. The disk radii are $1.898\,\mathrm{cm}$ for mass model A and
$1.900\,\mathrm{cm}$ for mass model B. These disks define the geometric region
used by the selection; they are not the complete physical windows or collimator
openings.

The remaining position uncertainty is handled explicitly. The catalogue gives
an analysis position for each hit, but does not resolve the interaction point
within the pixel. We represent that uncertainty by an axis-aligned box centred
on the recorded position, with half-lengths $0.075$, $0.075$, and
$0.150\,\mathrm{cm}$ along the analysis $x$, $y$, and $z$ axes. Its eight
vertices and centre give nine representative positions, not nine additional
measurements. For an assumed order of the first two hits, the two boxes produce
$9\times9=81$ position pairs and hence 81 candidate cones. Each cone is sampled
at 24 uniformly spaced azimuths and extended forward to the entrance plane. The
cone intersects the disk if a sampled point lies inside it or if a segment
joining adjacent samples crosses it. The test is deliberately permissive: one
intersection among the 81 candidate cones is enough to retain that order.
Figure~\ref{fig:compton_aperture_consistency} summarizes this construction.

\begin{figure}[H]
\centering
\includegraphics[width=\linewidth]{figures/manuscript/fig_compton_aperture_consistency_en.pdf}
\caption{Schematic of the Compton-cone/aperture-consistency veto. (a) The first
two TES hits define a reconstruction cone from their measured energies and
representative positions; the cone is a set of allowed source directions, not a
unique photon trajectory. (b) Each hit is represented by the eight vertices and
centre of an axis-aligned analysis box, giving $9\times9=81$ position hypotheses
for a specified order. (c) The cone trace in the entrance plane is approximated
by 24 azimuthal samples and their connecting segments. An intersection with the
analysis-aperture disk retains the event, whereas a valid reconstruction that
misses the disk is vetoed. Events without a valid cone and events with more than
six hits remain in the rate accounting as unreconstructed. The schematic is not
to scale.}
\label{fig:compton_aperture_consistency}
\end{figure}

The interaction order is treated separately from the within-pixel position
uncertainty. For a two-hit event, both orders are tested and either compatible
order is sufficient. For an event with three to six hits, all orders are first
enumerated using the recorded analysis positions. For a candidate sequence of
$n$ hits $\{(\mathbf r_j,\epsilon_j)\}_{j=1}^{n}$, the Compton sequence residual
is
\begin{equation}
 Q=\sum_{i=2}^{n-1}\left(\cos\theta_{C,i}-
 \hat{\mathbf u}_{i-1,i}\!\cdot\!\hat{\mathbf u}_{i,i+1}\right)^2,
\end{equation}
where $\theta_{C,i}$ is calculated from the energy $\epsilon_i$ deposited at
hit $i$ and the remaining energy $\sum_{k=i+1}^{n}\epsilon_k$; the two unit
vectors describe the adjacent flight segments. Orders without a physical
Compton solution do not enter the ranking. We select the order with the smallest
$Q$; if two orders are tied, we choose the one with the longer first lever arm.
Only the first-scatter cone of this selected order receives the 81-position
aperture test.

Brute-force enumeration is limited to $n\leq6$ for computational reasons. The
number of candidate orders grows as $n!$: a six-hit event requires at most
$6!=720$ orders, whereas seven and eight hits would require 5040 and 40320.
Each order entails $n-1$ kinematic checks and $n-2$ residual terms. After the
ranking, the fixed $9^2\times24=1944$ position--azimuth samples are evaluated
only for the selected minimum-$Q$ order. The six-hit limit therefore bounds the
factorial order search; it is not a detector or physics threshold.

The final decision is conservative. A reconstructable multi-hit event is vetoed
only when the cone or cones tested under the applicable ordering rule miss the
analysis aperture. An event with more than six hits, or with no physically valid
order, provides no usable aperture test under this implementation. It is marked
as unreconstructed and retained in the final rate rather than rejected. In this
way, failure to reconstruct an event is not itself treated as evidence that the
event is background.

\subsection{Mission-time trajectory and analytic rate folding}
\label{sec:mission_time_fold}

MeV balloon observations are background dominated: atmospheric secondaries and
cosmic-ray-induced activation set the rate floor, and both respond to the flight
environment \cite{Gallego2025Balloon,Kierans2017}. Detector-coupled Monte Carlo
through a fixed geometry provides the selected day-15 signal and background
rates under a common response, veto, and energy window, but those rates describe
only one longitude, latitude, and altitude. During a multi-day float the residual
atmospheric depth, geomagnetic cutoff, and activation inventory evolve, so a
constant-rate extrapolation of the day-15 snapshot is not a substitute for a
trajectory fold. We therefore map the reference rates onto a continuous 20 d
performance estimate and use five common-time anchors to determine the
accidental-coincidence correction to the background rate and the focused-signal
survival fraction.

The adopted trajectory is a \emph{synthetic reference profile}, not balloon
telemetry and not a source-visibility schedule. It contains 81 nodes from day 0
to day 20 at 0.25 d spacing, with a mild mid-latitude float envelope around
38.75 km altitude, $34^\circ$ latitude, and $100^\circ$ longitude. Altitude is
retained because changes in residual atmospheric depth affect both the particle
spectra that drive prompt background and activation and the atmospheric
transmission at 511\,keV multiplying the focused signal. Latitude and longitude supply a
controlled reference modulation of the geomagnetic environment; they do not
encode Earth occultation, repointing losses, or the visibility of a named source.
The fold is therefore a reference-exposure sensitivity calculation rather than a
forecast for a particular launch campaign.

We use the following indices throughout the fold: $g$ denotes the mass model,
$j$ the selection stage, $a$ an incident-particle family, $k$ a source-parent
nuclide, $r$ a common-time anchor, and $\ell$ an 81-node trajectory index. This
notation keeps the environmental scaling, detector response, and time grouping
distinct.

The selected monochromatic point-source rate before cross-event coincidence is
\begin{equation}
 S^{\mathrm{dir}}_g(t)=F_{\mathrm{TOA},511}A_{\mathrm{eff},g}T_{\mathrm{atm},511}(t),
 \label{eq:mission_signal}
\end{equation}
where $A_{\mathrm{eff},g}$ contains the detector response and the single-event
acceptance of the active-veto and topology/aperture selections. For an elevation $\epsilon$, the plane-parallel
slant transmission is
$T_{\mathrm{atm},511}(t)=\exp[-\mu_{\mathrm{air}}X(t)/\sin\epsilon]$; here
$\epsilon=45^\circ$ and $T_{15,45}=0.652034$. The reference significance curve
sets $F_{\mathrm{TOA},511}=\ftoaref$; $\fpayloadref=\ftoaref T_{15,45}$ is therefore the
same source normalization expressed at the day-15 payload plane, not a second
factor in Equation~\eqref{eq:mission_signal}.

The selected prompt rate is evaluated family by family,
\begin{equation}
 R_{\mathrm{prompt},g,j}(t)=\sum_a\Phi_a(t)K_{g,j,a},
 \label{eq:prompt_source_response}
\end{equation}
where each of the eight external fluxes $\Phi_a(t)$ has its own 81-node scale
and $K_{g,j,a}$ is the response coefficient of the geometry-specific transport
catalogue. The extrapolation changes the total scale of each family at a node;
its energy and angular templates remain those of the reference transport. It is
therefore an analytic reuse of detector-response templates, not 81 independent
environment-specific transports.

For activation, the production rate of source-parent $k$ from family $a$ is
\begin{equation}
 P_{g,a,k}(t)=\Phi_a(t)Y_{g,a,k},
 \label{eq:activation_production}
\end{equation}
where $Y_{g,a,k}$ is fixed by the geometry-specific production catalogue. Between
adjacent nodes the production rate is linearly interpolated, and the inventory
is advanced with the exact decay convolution,
\begin{equation}
 \begin{aligned}
 N_{g,a,k}(t_{\ell+1})={}&N_{g,a,k}(t_\ell)e^{-\lambda_k\Delta t_\ell}\\
 &+\int_0^{\Delta t_\ell}P_{g,a,k}(t_\ell+u)
 e^{-\lambda_k(\Delta t_\ell-u)}\,\mathrm du,\\
 A_{g,a,k}(t_\ell)={}&\lambda_kN_{g,a,k}(t_\ell).
 \end{aligned}
 \label{eq:delayed_inventory}
\end{equation}
The initial condition is $N_{g,a,k}(0)=0$. Thus the calculation contains neither
ground pre-activation nor inventory accumulated during ascent before the first
node; early threshold-crossing times are conditional on this boundary condition.
The selected delayed rate is then
\begin{equation}
 R_{\mathrm{delayed},g,j}(t)=
 \sum_a\sum_k A_{g,a,k}(t)\varepsilon_{g,j,a,k}.
 \label{eq:delayed_selected}
\end{equation}
This order of operations preserves the different half-lives and selected-decay
responses of the retained family--parent components. A single scale applied to
the total day-15 delayed rate would erase that information.

Let $B^{\mathrm{dir}}_g(t)$ be the sum of the selected seven-family prompt,
broadband atmospheric $\gamma$, and delayed rates before cross-event grouping. Common-time calculations are
performed at $r=1,\ldots,5$, corresponding to days 0, 5, 10, 15, and 20. For
anchor $r$, let $T^{\mathrm{rep}}_{g,r}$ be the equivalent statistical exposure,
$n^{\mathrm{tl}}_{g,r}$ the number of fully selected coincidence groups, and
$N^{\mathrm{probe}}_{g,r}$ and $n^{\mathrm{pass}}_{g,r}$ the inserted and
surviving focused-signal probes. We define two separate factors,
\begin{equation}
 \widehat\rho_{g,r}=
 \frac{n^{\mathrm{tl}}_{g,r}/T^{\mathrm{rep}}_{g,r}}
 {B^{\mathrm{dir}}_g(t_r)},\qquad
 \widehat\eta_{g,r}=
 \frac{n^{\mathrm{pass}}_{g,r}}{N^{\mathrm{probe}}_{g,r}}.
 \label{eq:anchor_time_factors}
\end{equation}
The background factor $\widehat\rho$ converts a direct component sum to a
common-time rate and can be slightly larger than unity when accidental grouping
moves events into the science window; it is not a survival probability.
$\widehat\eta$ is the conditional survival fraction of a weak focused event.
Using separate $\rho$ and $\eta$ factors avoids imposing a common
$e^{-R\tau}$ approximation on physically different quantities. The two factors
are interpolated independently and linearly between adjacent anchors. The
equivalent exposure is $2\times10^4$ s per anchor for mass model A and
$2\times10^5$ s per anchor for mass model B; these values set the sampling
support, not the 6 h physical step of the trajectory.

With $\Delta t_\ell=t_\ell-t_{\ell-1}$, cumulative counts use trapezoidal
integration,
\begin{equation}
 \begin{aligned}
 N_s(T_n)&=\sum_{\ell=1}^{n}\frac{S^{\mathrm{dir}}_g(t_{\ell-1})\eta_g(t_{\ell-1})+S^{\mathrm{dir}}_g(t_\ell)\eta_g(t_\ell)}{2}\,\Delta t_\ell,\\
 N_b(T_n)&=\sum_{\ell=1}^{n}\frac{B^{\mathrm{dir}}_g(t_{\ell-1})\rho_g(t_{\ell-1})+B^{\mathrm{dir}}_g(t_\ell)\rho_g(t_\ell)}{2}\,\Delta t_\ell.
 \end{aligned}
 \label{eq:cumulative_counts}
\end{equation}
We report $Z_G=N_s/\sqrt{N_b}$ and the Asimov measure
\begin{equation}
 Z_A=\sqrt{2\left[(N_s+N_b)\ln\left(1+\frac{N_s}{N_b}\right)-N_s\right]},
\end{equation}
and solve the same signal kernel for the $3\sigma$ and $5\sigma$ minimum
resolvable line fluxes.

Finite-transport uncertainty retains the correlation introduced by reuse of a
transport template: the node-dependent contribution of each retained event is
first integrated across all 81 nodes and only then squared. Independent
common-time-anchor, selected-area, and signal-probe terms are added afterward.
Threshold-crossing times are interpolated only between adjacent trajectory
nodes, and flux uncertainties are propagated with the delta method. These terms
describe finite Monte Carlo and common-time support; radiation-field,
cross-section, calibration, and pointing systematics are outside them. The
independent prompt-fold test and the five-anchor implementation diagnostic are
reported in Section~\ref{subsec:mission_fold_validation}.

\subsection{Exploratory response-weighted screening for non-balloon environments}
\label{subsec:environment_transfer_method}

We used an exploratory transfer calculation to ask which residuals of mass
model B would be most exposed to a change of radiation environment. The target
source fields are a 530\,km near-equatorial LEO proxy outside the South
Atlantic Anomaly (SAA), an exposed lunar-surface proxy, and quiet Sun--Earth L2
at two solar-modulation endpoints. The LEO source field follows the spectral
components used in low-Earth-orbit gamma-ray background calculations
\cite{Cumani2019LEO}. The lunar field combines the upward REDMoon spectrum for
Apollo-17 regolith at the 2019 solar minimum with the cosmic-photon sky visible
above the surface \cite{Dobynde2021REDMoon}. The Sun--Earth L2 source set
contains $\gamma$, proton, $\alpha$, electron, and positron components. Its
proton spectra use force-field modulation potentials of
$0.3793$\,GV for the 2009 solar-minimum/GCR-high endpoint and $0.803$\,GV for
the 2014 solar-maximum/GCR-low endpoint; the $\alpha$ component uses the
corresponding long-term GCR parameterization
\cite{Lotti2021Athena,Usoskin2005,Kuznetsov2017}. Energy denotes the total
kinetic energy of one incident particle, including the complete $\alpha$
nucleus.

The retained mass-model-B transfer response traces each selected prompt candidate to
its incident family and primary energy. A delayed candidate is instead coupled
to the weighted production-energy distribution with the same incident family,
parent-nuclide state, and production volume. If $e=0$ denotes the 38\,km
balloon reference, the represented rate for response stream
$x\in\{\gamma,\mathrm d\}$ is
\begin{equation}
 \widetilde B_{x,e}=
 \sum_{i\in\mathcal R_x}w_{i,0}
 \int p_i(E)
 \frac{\phi_{a_i,e}(E)}{\phi_{a_i,0}(E)}
 \,\mathrm dE ,
 \label{eq:environment_reweight}
\end{equation}
where $w_{i,0}$ is the selected balloon weight, $a_i$ is the incident family,
$\phi_{a,e}$ is the angle-domain-integrated source spectrum, and $p_i(E)$ is
the normalized primary-energy kernel. For a prompt $\gamma$ record, $p_i(E)$
is concentrated at its transported primary energy; for a delayed record it is
the corresponding activation-production distribution. Contributions are
retained only where the target and reference spectra have common non-zero
support. The direction template of the balloon response is not changed, so
Equation~\eqref{eq:environment_reweight} is an importance-reweighting
screening, not a new directional transport.

To place the environmental transfer and the mass-model-B background budget in
Section~\ref{sec:results} on the same normalization, we anchor the two
represented streams to the day-15 central rates,
\[
 B_{\gamma,0}=5.355349\times10^{-3}\cps,\qquad
 B_{\mathrm d,0}=3.580538\times10^{-3}\cps .
\]
With
\[
 c_\gamma=\frac{B_{\gamma,0}}{\widetilde B_{\gamma,0}}=0.924290,
 \qquad
 c_{\mathrm d}=\frac{B_{\mathrm d,0}}{\widetilde B_{\mathrm d,0}}=0.987825,
\]
we define the represented two-stream background and its environmental ratio as
\begin{equation}
 B^{(2)}_e=c_\gamma\widetilde B_{\gamma,e}
          +c_{\mathrm d}\widetilde B_{\mathrm d,e},
 \qquad
 {\cal R}^{(2)}_e=\frac{B^{(2)}_e}{B^{(2)}_0}.
 \label{eq:environment_two_stream}
\end{equation}
This is a stream-level re-anchoring approximation: it closes the current
stream totals but does not assert that every current event has the same
within-stream weight as the retained transfer sample.

The current seven-family prompt residual is
$1.811307\times10^{-3}\cps$, or 16.854\% of the full balloon rate, and is
supported by three positron records. The compact transfer response contains no
prompt-positron primary-energy kernel, so a family-resolved target rate cannot
be inferred for this term. To define a dimensionless index anchored to the
formal full-background balloon threshold, we make one explicit closure
convention: the fractional change of the complete background is represented by
$ {\cal R}^{(2)}_e$. Equivalently, the unmapped prompt share is carried with the
aggregate two-stream ratio only in Equation~\eqref{eq:environment_flux_screen};
it is not reported as an independently transferred positron rate. This
convention is the leading unmapped systematic of the screening index, which is
therefore not a complete target-background prediction.

For a common 20\,d exposure and selected effective area, the ratio of the
formal balloon signal kernel to an unattenuated, continuously on-source kernel
is
\begin{equation}
 \tau_{\mathrm B}=
 \frac{K_{\mathrm B}^{20\mathrm d}}
 {A_{\mathrm{eff,B}}T_{\mathrm{exp}}}=0.643861 .
 \label{eq:environment_signal_factor}
\end{equation}
We set $\tau_e=1$ for the non-balloon source planes and anchor the screening
index to the formal balloon Gaussian threshold
$F_{3\sigma,\mathrm B}=2.435681\times10^{-5}\phcms$:
\begin{equation}
 \frac{F_{3\sigma,e}^{\mathrm{screen}}}{F_{3\sigma,\mathrm B}}=
 \begin{cases}
  1, & e=0,\\[3pt]
  (\tau_{\mathrm B}/\tau_e)\sqrt{{\cal R}^{(2)}_e}, & e\ne0 .
 \end{cases}
 \label{eq:environment_flux_screen}
\end{equation}
The balloon anchor is a top-of-atmosphere flux; non-balloon values refer to the
local incident plane. The weight-concentration diagnostic
$N_{\mathrm{eff},w}=(\sum_iw_i)^2/\sum_iw_i^2$ is reported separately from the
formal total-background effective sample size. No target-specific activation
history, common-time veto or accidental-survival factor, visibility schedule,
or platform geometry enters this calculation.

\section{Results}
\label{sec:results}

The results follow the sequence in Figure~\ref{fig:workflow}. We first establish the background and response of mass model A and trace the selected events to the materials, parent nuclides, and spatial regions that limit its performance. We then test the complete side-chimney configuration of mass model B and compare the two designs on the common mission time axis. The energy response, science window, arrival-time grouping, and Compton topology/aperture criterion are the same for both models; the active-shield channels are part of each complete design.

\subsection{Baseline background and selection performance of mass model A}
\label{subsec:model_a_baseline}
\label{subsec:reference_background_results}

Mass model A defines the reference performance. Table~\ref{tab:phase2_cutflow} gives the day-15 direct, no-coincidence cut flow in the $[510.58,511.42)\keV$ science window. The joint BGO and plastic-scintillator gate reduces the total background from 3,086 records and $5.43630\cps$ to 443 records and $6.07099\times10^{-2}\cps$, removing 98.883\% of the pre-veto rate. The seven prompt particle families fall from $2.91217\cps$ to a zero central value, the broadband atmospheric $\gamma$ component falls from $2.35667\cps$ to $1.78923\times10^{-2}\cps$, and delayed activation falls from $0.167460\cps$ to $4.28176\times10^{-2}\cps$.

After the Compton topology/aperture selection, 400 background records remain at $5.30571\times10^{-2}\cps$, or 87.394\% of the active-veto rate. All 28,612 isolated focused-signal records in the measured window pass the active gate, and 28,006 survive the final selection. The corresponding effective area is $15.12324\pm0.04492\,\mathrm{cm^2}$, equal to 97.882\% of the measured-window signal area. Active anticoincidence supplies most of the rate reduction, while the topology selection removes a smaller residual and retains most of the focused signal.

\begin{table}[H]
\centering
\caption{Mass model A science-window cut flow. Background entries are day-15 direct, no-coincidence record counts and rates; signal entries give record counts and effective area for the common 37,194-photon focal-plane sample.}
\label{tab:phase2_cutflow}
\footnotesize
\begin{tabular}{lccc}
\toprule
Stream & Pre-veto & After active veto & Final topology/aperture \\
\midrule
Seven prompt particle families & 869 / $2.91217\cps$ & 0 / 0 & 0 / 0 \\
Broadband atmospheric $\gamma$ & 922 / $2.35667\cps$ & 7 / $1.78923\times10^{-2}\cps$ & 6 / $1.53363\times10^{-2}\cps$ \\
Delayed activation & 1295 / $0.167460\cps$ & 436 / $4.28176\times10^{-2}\cps$ & 394 / $3.77209\times10^{-2}\cps$ \\
Total background & 3086 / $5.43630\cps$ & 443 / $6.07099\times10^{-2}\cps$ & 400 / $5.30571\times10^{-2}\cps$ \\
Focused signal & 28612 / $15.45048\,\mathrm{cm^2}$ & 28612 / $15.45048\,\mathrm{cm^2}$ & 28006 / $15.12324\,\mathrm{cm^2}$ \\
\bottomrule
\end{tabular}
\end{table}

The nonzero components of the final direct background are delayed activation
and broadband atmospheric $\gamma$ rays
(Table~\ref{tab:reference_background_budget}). Delayed activation contributes
$3.77209\times10^{-2}\cps$ (71.095\%), and the broadband component contributes
$1.53363\times10^{-2}\cps$ (28.905\%). Because the event weights differ among
components, record count alone does not measure statistical support; the
400-record total has $N_{\mathrm{eff}}=46.77$.

\begin{table}[H]
\centering
\caption{Day-15 final direct background of mass model A. Standard errors and effective sample sizes are calculated from $\sum_i w_i^2$ for unequal event weights.}
\label{tab:reference_background_budget}
\footnotesize
\begin{tabular}{lrrrrr}
\toprule
Component & Records & Rate / $\cps$ & Standard error / $\cps$ & $N_{\mathrm{eff}}$ & Total / \% \\
\midrule
Delayed activation & 394 & $3.77209\times10^{-2}$ & $4.58152\times10^{-3}$ & 67.79 & 71.095 \\
Broadband atmospheric $\gamma$ & 6 & $1.53363\times10^{-2}$ & $6.26100\times10^{-3}$ & 6.00 & 28.905 \\
\midrule
Total & 400 & $5.30571\times10^{-2}$ & $7.75825\times10^{-3}$ & 46.77 & 100.000 \\
\bottomrule
\end{tabular}
\end{table}

\subsection{Background origins and spatial coupling in mass model A}
\label{subsec:mass_model_a_origins}

The total rate fixes the background scale but does not explain why events pass the final selection. We therefore traced the selected mass-model-A catalogue by initiating particle family, material in which the radioactive parent was produced, source-volume/parent-nuclide pair, and spatial region relative to the TES.

Proton-induced activation is the largest single contribution, with 35 records and $2.48463\times10^{-2}\cps$, or 46.829\% of the total direct background. The remaining delayed terms are $7.49239\times10^{-3}\cps$ from $\alpha$-induced activation, $3.60054\times10^{-3}\cps$ from neutron-induced activation, $1.68617\times10^{-3}\cps$ from $\gamma$-induced activation, $7.02083\times10^{-5}\cps$ from $e^+$-induced activation, and $2.52755\times10^{-5}\cps$ from $e^-$-induced activation. Together they reproduce the final delayed rate of $3.77209\times10^{-2}\cps$.

Copper accounts for 70.87\% of the weighted selected delayed rate. The leading identified source-volume/nuclide pairs are $^{62}$Cu in the L2 copper support ring, the 50 mK mixing-chamber copper plate, and the L0 TES heat-sink ring. They are followed by $^{199}$Po in the Bi upper half-cylinder and $^{64}$Cu in the mixing-chamber plate and L2 support ring. This order already includes transport, detector response, and final event selection; it cannot be inferred from material mass or total activity alone.

The spatial decomposition gives the corresponding geometric coupling. The DR mixing chamber and staged cold plates supply 32.18\% of the weighted selected delayed rate, and TES-near structures supply 47.29\%, for a combined share of 79.47\%. Their rates are approximately $1.21398\times10^{-2}$ and $1.78397\times10^{-2}\cps$. The limiting feature is therefore not the total copper inventory, but the coupling of activated copper and nearby structures to events that satisfy the TES selection.

This diagnosis defines the design test addressed by mass model B: whether changing the relative position of the focal plane and cold structures, while surrounding the detector with local active shielding, can reduce the selected coupling without losing focused response. Because the entrance, local tungsten structure, BGO layout, and part of the nearby material inventory change together, the comparison tests the complete configuration rather than the chimney displacement alone.

\FloatBarrier

\subsection{From the source diagnosis to the complete mass model B configuration}
\label{subsec:optimized_shield_results}

Mass model B places the six TES layers along the lateral chimney axis at $z'=-2.80$ cm. The aluminium chimney has inner and outer radii of 10.75 and 11.05 cm, respectively, and a 3 mm wall joined continuously to the dilution-refrigerator shell; its bottom is 0.25 cm above the DR bottom. A four-sided tungsten frame with 6.00 and 5.40 cm outer and inner square openings and a 2.00 cm axial depth supports the focal plane. Three active BGO volumes cover the chimney side, the front optical annulus, and the rear cold-port annulus. Mass model B has no plastic-scintillator channel, and the active criterion requires the summed BGO deposit in each $1\,\mu\mathrm{s}$ group to remain below 50 keV.

Figure~\ref{fig:optimized_shield_background} compares the two cold-end
structures in the same instrument coordinates and at the same scale, and
Table~\ref{tab:final_geometry} lists the model-B dimensions used in transport.
The detector-neighbourhood panels show the TES, key copper supports, tungsten
aperture/frame, and active BGO. After final selection, the focused effective
areas are $15.12324\pm0.04492\,\mathrm{cm^2}$ for mass model A and
$15.08544\pm0.04503\,\mathrm{cm^2}$ for mass model B; model B retains 99.75\%
of the model-A area.

\begin{figure}[!ht]
\centering
\includegraphics[width=\linewidth]{figures/section4_current/fig09_geometry_response_detailed.pdf}
\caption{Dimensionally aligned $x'$--$z'$ sections of the cold end and focal plane, overlaid with the source positions of final selected delayed events. Panels (a) and (b) show mass models A and B over identical coordinate ranges and at the same scale; panels (c) and (d) show their detector neighbourhoods at a common physical scale. The drawings identify the staged copper cold plates, six-layer TES, active BGO, tungsten aperture/frame, and the model-B chimney and copper cold finger. Each symbol marks the production position associated with a day-15 delayed event that survives the final selection; colour identifies the source material, and marker area follows the selected weighted rate on a common scale subject to a minimum display size. Sections are drawn to the transport-geometry dimensions; minor CSG apertures and service components are omitted for clarity.}
\label{fig:optimized_shield_background}
\end{figure}

\begin{table}[!ht]
\centering
\caption{Mass model B side-chimney geometry in local instrument coordinates.}
\label{tab:final_geometry}
\footnotesize
\begin{tabular}{p{0.23\linewidth}p{0.34\linewidth}p{0.34\linewidth}}
\toprule
Element & Final setting & Design role \\
\midrule
Aluminium chimney & inner/outer radii 10.75/11.05 cm; 3 mm wall & places the TES laterally outside the staged cold region \\
Common axis and clearance & $z'=-2.80$ cm; chimney/DR bottoms $-13.85/-14.10$ cm & continuous interface with 0.25 cm clearance \\
TES focal plane & six Ta-pixel layers along the chimney axis & preserves detector response and multilayer stopping depth \\
Tungsten focal frame & outer/inner square 6.00/5.40 cm; depth 2.00 cm & local support and aperture definition \\
Cold/optical ports & 1.50 cm cold port; 6.00 cm upper optical opening & retains thermal services and focused-beam access \\
Active BGO & side shield, front optical annulus, rear cold-port annulus & anticoincidence coverage around the chimney \\
BGO threshold & 50 keV offline deposited energy & active-veto decision \\
\bottomrule
\end{tabular}
\end{table}

The agreement in effective area constrains the focused-photon sample; it does not imply equal acceptance for diffuse photons. Mass model A uses an entrance envelope of radius approximately 1.898 cm and a 624-hole tungsten collimator, whereas mass model B uses a 2.70 cm-radius physical window and an open four-sided tungsten frame. The active-shield configurations also differ. The reported model-B background and sensitivity are therefore conditional on its entrance, tungsten frame, BGO layout, and active-channel configuration; no single-factor ablation is available.

\FloatBarrier

\subsection{Full-chain background and source redistribution in mass model B}
\label{subsec:optimized_fullchain}

Mass model B uses the same 420 eV FWHM TES response, $[510.58,511.42)\keV$ science window, and Compton topology/aperture test as mass model A. Table~\ref{tab:optimized_cutflow} gives its day-15 direct, no-coincidence cut flow.

Active BGO reduces the total background from 7,684 records and $2.38372\cps$ to 137 records and $1.11379\times10^{-2}\cps$, or 0.4672\% of the pre-veto rate. All 28,434 isolated focused-signal records in the measured window pass the active gate. The final Compton topology/aperture selection leaves 126 background records and retains 96.492\% of the active-veto weighted rate. It retains 27,936 focused records, corresponding to $15.08544\pm0.04503\,\mathrm{cm^2}$ and 98.249\% of the measured-window signal area.

\begin{table}[H]
\centering
\caption{Mass model B response-convolved science-window cut flow. Background entries are day-15 direct, no-coincidence record counts and rates; signal entries give record counts and effective area.}
\label{tab:optimized_cutflow}
\footnotesize
\begin{tabular}{lccc}
\toprule
Stream & Pre-veto & After active BGO & Final topology/aperture \\
\midrule
Seven prompt particle families & 1448 / $0.878692\cps$ & 3 / $1.81131\times10^{-3}\cps$ & 3 / $1.81131\times10^{-3}\cps$ \\
Broadband atmospheric $\gamma$ & 2987 / $1.37944\cps$ & 12 / $5.35535\times10^{-3}\cps$ & 12 / $5.35535\times10^{-3}\cps$ \\
Delayed activation & 3249 / $0.125586\cps$ & 122 / $3.97120\times10^{-3}\cps$ & 111 / $3.58054\times10^{-3}\cps$ \\
Total background & 7684 / $2.38372\cps$ & 137 / $1.11379\times10^{-2}\cps$ & 126 / $1.07472\times10^{-2}\cps$ \\
Focused signal & 28434 / $15.35436\,\mathrm{cm^2}$ & 28434 / $15.35436\,\mathrm{cm^2}$ & 27936 / $15.08544\,\mathrm{cm^2}$ \\
\bottomrule
\end{tabular}
\end{table}

The final model-B direct background contains $1.81131\times10^{-3}\cps$ from
the seven prompt particle families, $5.35535\times10^{-3}\cps$ from broadband
atmospheric $\gamma$ rays, and $3.58054\times10^{-3}\cps$ from delayed
activation. Their shares are 16.854\%, 49.830\%, and 33.316\%, respectively.
Relative to mass model A, the delayed and broadband-$\gamma$ rates fall to
9.492\% and 34.920\%, and the total rate falls to 20.256\%, a factor of 4.94.

The component balance changes with the complete configuration. Mass model A is
activation dominated, whereas the broadband atmospheric $\gamma$ component is
the largest component in mass model B. Table~\ref{tab:optimized_component_precision}
shows the unequal-weight support of this final composition. The nonzero prompt
central value in model B is supported by three positron-family templates and
has a broad finite-sample interval; it enters the total-rate calculation but
does not establish a stable positron-dominated component. Likewise, the zero
prompt central value in mass model A does not imply that prompt background is
absent.

\begin{table}[!ht]
\centering
\caption{Finite-transport support of the day-15 final direct background in mass model B. Standard errors and effective sample sizes are calculated from $\sum_i w_i^2$ for unequal event weights.}
\label{tab:optimized_component_precision}
\scriptsize
\resizebox{\linewidth}{!}{%
\begin{tabular}{lrrrrr}
\toprule
Component & Records & Rate / $\cps$ & Standard error / $\cps$ & $N_{\mathrm{eff}}$ & Total / \% \\
\midrule
Seven prompt particle families & 3 & $1.81131\times10^{-3}$ & $1.04576\times10^{-3}$ & 3.00 & 16.854 \\
Broadband atmospheric $\gamma$ & 12 & $5.35535\times10^{-3}$ & $1.55631\times10^{-3}$ & 11.84 & 49.830 \\
Delayed activation & 111 & $3.58054\times10^{-3}$ & $5.36971\times10^{-4}$ & 44.46 & 33.316 \\
\midrule
Total & 126 & $1.07472\times10^{-2}$ & $1.95040\times10^{-3}$ & 30.36 & 100.000 \\
\bottomrule
\end{tabular}}
\end{table}

Copper-derived events contribute $(2.26708\pm0.42826)\times10^{-3}\cps$, or
63.32\% of the final model-B delayed rate. Figure~\ref{fig:geometry_response_rates}
places this material projection beside the independent spatial projection. The
DR/mixing-chamber and staged-cold contribution changes from
$(1.21398\pm0.26869)\times10^{-2}\cps$ in mass model A to
$(1.40604\pm1.09289)\times10^{-4}\cps$ in mass model B. The latter estimate is
supported by only seven source records and has a large relative error. The
total delayed rate changes from $(3.77209\pm0.45815)\times10^{-2}$ to
$(3.58054\pm0.53697)\times10^{-3}\cps$. Identified model-B groups still include
$^{62}$Cu in the L5 copper ring and the TES heat sink; the selected copper
contribution is reduced and spatially redistributed rather than eliminated.

\begin{figure}[!ht]
\centering
\includegraphics[width=\linewidth]{figures/section4_current/fig09b_geometry_response_rates.pdf}
\caption{Selected delayed-background redistribution and complete-design response ratios. (a) Day-15 final delayed detector rate from the cold stages, TES-near structures, and other structures in mass models A and B. (b) Model-B/model-A ratios for the cold-stage, Cu-derived, and total delayed components, the total science-window background, and the final selected signal effective area; the vertical dashed line denotes equality. Error bars are finite-transport $1\sigma$ uncertainties from the unequal-weight second moment, with the two models propagated independently in each ratio. They do not include common-time or physical systematics. Because the entrance geometry, local tungsten structure, BGO arrangement, nearby material inventory, and active-channel configuration differ, panel (b) compares complete designs rather than a single-parameter ablation.}
\label{fig:geometry_response_rates}
\end{figure}

Figure~\ref{fig:day15_selection_spectra} shows the day-15 spectra before active
anticoincidence, after the active gate, and after the final Compton selection.
The 480--550 keV band provides a selection diagnostic independent of the
narrow-window integral. In mass model A, its weighted rate changes from
$10.7984\cps$ to $0.140797\cps$ and then $0.121285\cps$; the corresponding
model-B rates are $4.92157$, $0.0735321$, and $0.0655197\cps$. The same ordering
appears in the science-window cut flows: active anticoincidence supplies the
dominant suppression, and the Compton topology/aperture selection makes a
smaller correction to the residual.

\begin{figure}[!ht]
\centering
\includegraphics[width=\linewidth]{figures/section4_current/fig05_day15_selection_spectra_current.pdf}
\caption{Day-15 detector-side differential count-rate spectra through the event-selection chain. Panels (a) and (b) show 480--550 keV for mass models A and B; four adjacent 0.25 keV analysis bins are summed into each 1 keV display bin. Panels (c) and (d) show 508--514 keV at the native 0.25 keV binning; the shaded interval and dotted boundaries mark the $[510.58,511.42)\keV$ science window. Curves show the direct, no-coincidence background before active anticoincidence, after active anticoincidence, and after the final Compton topology/aperture selection. The translucent band in the broad-band panels and the bin-wise error bars in the zoom panels show finite-transport $1\sigma$ uncertainties from $\sqrt{\sum_i w_i^2}$. Zero-weight bins are left blank, without smoothing or pseudocounts. The ordinate is the detector differential count rate after transport and TES-response convolution.}
\label{fig:day15_selection_spectra}
\end{figure}

The detector-side results connect the reference rate, the source diagnosis, the
candidate geometry, and the repeated selection. Both delayed activation and
the broadband atmospheric $\gamma$ component have lower central rates in mass
model B than in mass model A, and the broadband component is the largest
model-B residual. This supports the conditional performance assessment of mass
model B as a complete candidate design, but it does not isolate the lateral
chimney as a single causal factor.

\FloatBarrier

\subsection{Common-time fold and mission performance}
\label{subsec:mission_fold_validation}
\label{subsec:final_mission_performance}

Before the mission projection, the family-resolved prompt environmental fold was tested by direct transport for three independent off-reference source definitions. The three direct-to-analytic rate ratios are 0.992, 1.004, and 0.994, with a maximum difference of $0.51\sigma$. The test covers the $e^+$, neutron, and $\gamma$ prompt slices in the 480--550 keV band.

At the day-0, 5, 10, 15, and 20 anchors, Poisson arrivals are regenerated from the physical component rates and passed through the common time grouping, detector response, active anticoincidence, and Compton selection. The ratio of the common-time final rate to the direct rate defines $\widehat\rho_g$, and focused-signal replay gives the accidental survival $\widehat\eta_g$; both are interpolated to the 81 mission nodes. The model-A $\widehat\rho_g$ values span 0.9307--1.0060, and the model-B values span 0.9901--1.0293 (Figure~\ref{fig:common_time_anchors}).

\begin{figure}[H]
\centering
\includegraphics[width=\linewidth]{figures/section4_current/fig04_common_time_normalization.pdf}
\caption{Five-anchor common-time diagnostic. (a) Final direct and common-time-corrected science-window rates for mass models A and B; error bars are Poisson standard errors at fixed exposure. (b) Background correction $\rho_g(t)$ and focused-signal survival $\eta_g(t)$; the latter error bars are binomial. The equivalent exposure per anchor is $2\times10^4$ s for mass model A and $2\times10^5$ s for mass model B. All anchors use the same $1\,\mu\mathrm{s}$ grouping rule, and the curves are linear interpolants of the five anchors.}
\label{fig:common_time_anchors}
\end{figure}

The focused-signal $\widehat\eta_g$ values span 0.96810--0.97039 for mass model A and 0.99547--0.99580 for mass model B; their day-15 values are 0.9682365 and 0.995472. The higher model-B survival follows from lower coincidence occupancy rather than a different signal selection. The interpolated $\rho_g$ and $\eta_g$ values are combined with the same trajectory transmission and component responses and integrated over 81 nodes at 0.25 d spacing.

At 20 d, mass model A accumulates 3928.83 signal counts at the reference flux
and 87632.34 background counts; mass model B accumulates 4028.15 and 18568.89.
The broadband atmospheric $\gamma$ and remaining-background totals are
25760.39 and 61871.95 counts for mass model A and 9359.44 and 9209.45 counts for
mass model B. The latter term combines the seven prompt particle families and
delayed activation and is not an activation-only quantity. The model-B
cumulative background is 21.19\% of the model-A value.

Figure~\ref{fig:optimized_mission_significance} and
Table~\ref{tab:primary_sensitivity} summarize the mission projection. At
$\ftoaref$, the 20 d Gaussian significances are 13.272 for mass model A and
29.561 for mass model B; the corresponding Asimov values are 13.174 and
28.578. The Gaussian $3\sigma$ minimum resolvable top-of-atmosphere fluxes are
$(5.425\pm0.403)\times10^{-5}$ and
$(2.436\pm0.223)\times10^{-5}\phcms$; the Asimov values are
$(5.434\pm0.403)\times10^{-5}$ and
$(2.445\pm0.223)\times10^{-5}\phcms$. The quoted errors are local $1\sigma$
statistical uncertainties propagated from the finite-transport estimates rather
than joint confidence intervals.

\begin{figure}[!ht]
\centering
\includegraphics[width=0.98\linewidth]{figures/section4_reference_flux_20260830/fig10_mission_performance.pdf}
\caption{Conditional 81-node mission performance of mass models A and B. (a)
Gaussian counting significance at the top-of-atmosphere reference flux
$\ftoaref=2.4\times10^{-4}\phcms$; horizontal lines mark $3\sigma$ and
$5\sigma$. (b) Gaussian and Asimov $3\sigma$ minimum resolvable
top-of-atmosphere flux. Both designs use the same source, detector response,
science window, timing, topology/aperture test, and atmospheric-transmission
conventions; their entrances, local tungsten structures, and geometry-specific
active channels differ.}
\label{fig:optimized_mission_significance}
\end{figure}

\begin{table}[!ht]
\centering
\caption{Conditional mission projection for mass models A and B. Flux-threshold errors are local $1\sigma$ statistical uncertainties propagated from the finite-transport estimates.}
\label{tab:primary_sensitivity}
\footnotesize
\resizebox{\linewidth}{!}{%
\begin{tabular}{lrr}
\toprule
Quantity & Mass model A & Mass model B \\
\midrule
Selected effective area / $\mathrm{cm^2}$ & $15.12324\pm0.04492$ & $15.08544\pm0.04503$ \\
Day-15 direct background / $\cps$ & $5.30571\times10^{-2}$ & $1.07472\times10^{-2}$ \\
Day-15 signal at $\ftoaref$ / $\cps$ & $2.29144\times10^{-3}$ & $2.35000\times10^{-3}$ \\
20 d background counts & 87632.34 & 18568.89 \\
20 d signal counts at $\ftoaref$ & 3928.83 & 4028.15 \\
$Z_{20\mathrm d}$, Gaussian / Asimov & 13.272 / 13.174 & 29.561 / 28.578 \\
$T_{3\sigma}$ at $\ftoaref$, Gaussian / Asimov / d & 0.669 / 0.684 & 0.198 / 0.206 \\
$T_{5\sigma}$ at $\ftoaref$, Gaussian / Asimov / d & 1.913 / 1.954 & 0.452 / 0.487 \\
$F_{3\sigma}(20\,\mathrm d)$, Gaussian / Asimov / $\phcms$ &
$(5.425\pm0.403)\times10^{-5}$ / $(5.434\pm0.403)\times10^{-5}$ &
$(2.436\pm0.223)\times10^{-5}$ / $(2.445\pm0.223)\times10^{-5}$ \\
$F_{5\sigma}(20\,\mathrm d)$, Gaussian / Asimov / $\phcms$ &
$(9.042\pm0.671)\times10^{-5}$ / $(9.067\pm0.671)\times10^{-5}$ &
$(4.059\pm0.372)\times10^{-5}$ / $(4.084\pm0.372)\times10^{-5}$ \\
\bottomrule
\end{tabular}}
\end{table}

\FloatBarrier

The Gaussian threshold ratio separates into cumulative-background and cumulative-signal-kernel terms. The 20 d model-B background is 21.19\% of the model-A value, while its cumulative signal kernel is 1.025278 times the model-A value:
\begin{equation}
 \frac{F_{\min,A}}{F_{\min,B}}
 =\sqrt{\frac{N_{b,A}}{N_{b,B}}}\frac{K_B^{20\mathrm d}}{K_A^{20\mathrm d}}
 =2.172397\times1.025278
 =2.227311 .
 \label{eq:fmin_ratio_decomposition}
\end{equation}
Here $K_g^{20\mathrm d}$ is the cumulative signal kernel integrated over the 81 nodes. Equation~\eqref{eq:fmin_ratio_decomposition} shows that most of the 2.227-fold threshold difference follows from the square root of the background ratio, with a smaller correction from the cumulative signal kernel.

Mass model B retains 99.75\% of the focused effective area of mass model A and reduces the 20 d cumulative background and Gaussian $3\sigma$ threshold to 21.19\% and 44.90\%, respectively. Because the entrance, tungsten frame, active shielding, and nearby materials change together, this is a conditional comparison of two complete designs and does not identify a single geometric cause.

\FloatBarrier

\section{Conditional transfer to orbital and lunar-surface environments}
\label{sec:environment_screening}

The selected mass-model-B events connect the narrow detector window to a much
wider range of primary energies. The response-weighted 10--90\% intervals are
$0.598$--$29.6$\,MeV for prompt broadband $\gamma$ rays and
$0.0745$--$611$\,GeV, $0.115$--$2.01$\,GeV, $1.37$--$48.4$\,GeV,
$18.7$--$82.9$\,GeV, and $0.0197$--$195$\,GeV for delayed $\gamma$,
positron, proton, $\alpha$, and neutron production, respectively. The
$\alpha$ interval is expressed in total kinetic energy per nucleus. These
family-dependent bands define the primary-energy support of the spectral
reweighting in Equation~\eqref{eq:environment_reweight}.

\begin{figure}[!ht]
\centering
\includegraphics[width=\linewidth]{figures/environment_screening_20260830/fig11_environment_transfer_en.pdf}
\caption{Response-weighted environmental screening for mass model B.
(a) Angle-domain-integrated $\gamma$-ray source spectra; the shaded
band is the 10--90\% primary-energy interval of the selected prompt-$\gamma$
response. (b) Proton spectra, including the 2009 GCR-high and 2014 GCR-low
Sun--Earth L2 endpoints; the shaded band is the corresponding delayed-response
interval. Shading indicates response support, not source uncertainty. Each
source retains its declared physical angular domain, so the curves are not
common-isotropic-intensity measurements. (c) The two-stream Gaussian
$3\sigma$ screening index from Equation~\eqref{eq:environment_flux_screen};
numeric labels are in $10^{-5}\phcms$. Open markers emphasize that the
non-balloon values are response-transfer indices, not matched-transport
sensitivities or confidence intervals.}
\label{fig:environment_transfer}
\end{figure}

Table~\ref{tab:environment_screening} gives the resulting central indices. The
quiet non-SAA LEO source field yields $ {\cal R}^{(2)}=0.5627$. Removal of the
balloon attenuation in the signal kernel then gives a threshold index of
0.4830, corresponding to $1.176\times10^{-5}\phcms$. This apparent advantage
is limited to the adopted quiet non-SAA proxy; it does not include trapped
particles or activation accumulated during and after SAA passages.

The exposed lunar proxy instead gives $ {\cal R}^{(2)}=23.02$, driven mainly by
the regolith-related $\gamma$ field in the prompt response band. The
unattenuated signal partly offsets that increase, leaving a screening index of
3.089 and $F_{3\sigma}^{\mathrm{screen}}=7.525\times10^{-5}\phcms$. The value
applies to the adopted exposed-surface field and isolated detector--cryostat
model, not to a lander or lunar-base configuration.

\begin{table}[!ht]
\centering
\caption{Two-stream environmental screening for mass model B. The represented
background contains the broadband-$\gamma$ and delayed-activation streams.
$F_{3\sigma}^{\mathrm{screen}}$ is normalized to the current formal balloon
Gaussian threshold through Equation~\eqref{eq:environment_flux_screen}; the
balloon value is quoted at the top of the atmosphere and the others at the
local incident plane. This normalization assumes that the fractional change of
the complete background follows the represented two-stream ratio; it does not
infer a separate target rate for the unmapped prompt-positron residual.
$N_{\mathrm{eff},w}$ describes concentration of the transfer weights and is
neither the formal total-background effective sample size nor a confidence
interval.}
\label{tab:environment_screening}
\scriptsize
\resizebox{\linewidth}{!}{%
\begin{tabular}{llrrrr}
\toprule
Scenario & Source-field proxy &
$ {\cal R}^{(2)}_e$ &
$F_{3\sigma,e}^{\mathrm{screen}}/F_{3\sigma,\mathrm B}$ &
$F_{3\sigma,e}^{\mathrm{screen}}$ / $\phcms$ &
$N_{\mathrm{eff},w}$ \\
\midrule
38 km balloon & formal 81-node trajectory &
1.0000 & 1.0000 & $2.4357\times10^{-5}$ & -- \\
530 km LEO & near-equatorial, quiet non-SAA source field &
0.5627 & 0.4830 & $1.1764\times10^{-5}$ & 44.27 \\
Exposed lunar surface & REDMoon Apollo-17 regolith proxy &
23.0228 & 3.0894 & $7.5247\times10^{-5}$ & 5.23 \\
Sun--Earth L2 (2014) & solar maximum, GCR-low &
5.9460 & 1.5700 & $3.8241\times10^{-5}$ & 5.73 \\
Sun--Earth L2 (2009) & solar minimum, GCR-high &
11.9177 & 2.2227 & $5.4139\times10^{-5}$ & 5.05 \\
\bottomrule
\end{tabular}}
\end{table}

At Sun--Earth L2, the ordering of the two endpoints follows the proton spectra:
stronger solar modulation in 2014 lowers the GCR-driven delayed response,
whereas the weakly modulated 2009 field raises it. The corresponding screening
range is $(3.824$--$5.414)\times10^{-5}\phcms$, or 1.570--2.223 times the
formal balloon threshold. The numerical amplitude is weakly supported,
however. Four selected $^{62}$Cu survivors share one activation key represented
by one 1.366\,GeV proton-production record, and
$N_{\mathrm{eff},w}=5.05$--5.73. The response therefore establishes a
physically interpretable proton-activation risk, not a converged L2
background.

The missing physics differs by environment. The LEO proxy omits
inclination- and time-dependent geomagnetic cutoffs, trapped and SAA particles,
post-SAA decay, Earth-limb history, spacecraft mass, attitude, and duty cycle.
The lunar proxy omits local regolith and terrain, the lander or base, power
systems, solar energetic particles, and heavy ions. The Sun--Earth L2 proxy
omits solar energetic particles, heavy ions, directional soft protons,
spacecraft structure and attitude, and orbit-specific interruptions. None of
the three target fields has a matched prompt--activation--delayed transport or
its own common-time event fold. These indices therefore identify follow-up
background drivers; they do not rank completed mission concepts.

\FloatBarrier

\section{Discussion}
\label{sec:discussion}

Selected-event provenance links the model-A background diagnosis to the
model-B performance comparison. In mass model A, the TES lies below the staged cold
plates; this region supplies 32.18\% of the weighted selected delayed science-window
rate, and copper accounts for 70.87\% in the independent material projection. In
mass model B, the six-layer focal plane is placed laterally in a BGO-surrounded
chimney. The selected cold-stage rate then changes from
$(1.21398\pm0.26869)\times10^{-2}$ to
$(1.40604\pm1.09289)\times10^{-4}\cps$, although the latter estimate is based on
only seven source records. The selected effective area changes by less than
0.25\%, while the day-15 direct background, 20 d cumulative background, and
selected delayed rate decrease by factors of 4.94, 4.72, and 10.53,
respectively. The Gaussian $3\sigma$ flux threshold improves by a factor of
2.227.

This analysis sequence is consistent with recent instrument-background studies.
COSI transports a detailed balloon payload, separates prompt and activation
components, and compares their spectral and temporal behaviour with flight data
\cite{Gallego2025Balloon}. DIXE likewise begins with a detailed TES payload and
asks which incident particles, production regions, and processes make the
detected spectrum \cite{Tian2026DIXE}. Here the distinction between total
inventory and selected-event origin is decisive: the leading selected groups are
localized copper and near-field structures, so designing from total activity
alone would miss the geometric lever.

This origin map also explains why simply adding passive material below the cold
copper is an inefficient remedy. A passive layer can attenuate a 511 keV photon
through Compton scattering and photoelectric absorption, but complete coverage
requires thickness and mass, competes with the optical and cryogenic openings,
and adds material that can itself be activated \cite{NISTXCOM}. A 511 keV
photon is below the $2m_ec^2=1.022$ MeV nuclear-field pair-production
threshold and therefore cannot itself create an $e^-e^+$ pair. Higher-energy
broadband photons and particle cascades can, however, produce positrons in added
material, and $\beta^+$ activity can do the same; their annihilation then
creates secondary 511 keV photons. The model-B layout places the TES outside the
direct projection of the activated copper and is intended to reduce this coupling
before relying on active shielding. The compact BGO layout of mass model B also points
to a possible active-shield mass saving, but that saving is not quantified by
the present complete-design comparison.

The selection spectra also separate the roles of the instrument layers. Active
BGO supplies the dominant suppression in mass model B, reducing the direct
science-window rate from $2.38372$ to $1.11379\times10^{-2}\cps$. The final
topology/aperture test leaves $1.07472\times10^{-2}\cps$ while retaining
98.249\% of the measured-window focused signal. The final rate consists of
49.830\% broadband atmospheric $\gamma$ rays, 33.316\% delayed activation, and 16.854\%
from the seven prompt particle families. The prompt central value is supported
by only three selected records, so this composition is more securely interpreted
as a change from activation dominance in model A to a mixed residual in model B
than as a precise ordering of all three components.

The observing role of this telescope is complementary to INTEGRAL/SPI and COSI,
which measure diffuse bulge, disk, and broader Galactic emission
\cite{Siegert2016AA,Kierans2020,Siegert2020COSI,Yoneda2025,Karwin2023}. The
Laue-lens TES concept exchanges sky coverage for a concentrated beam and a
sub-keV window. Its natural targets are compact 511 keV sources, an unresolved
central component, and pointed follow-up observations. The 20 d values in
Table~\ref{tab:primary_sensitivity} quantify that continuously observed,
narrow-line reference case.

\subsection{Cross-environment interpretation}

The screening in Section~\ref{sec:environment_screening} does not alter the
formal balloon result. Its robust content is the separation of environmental
drivers: the quiet non-SAA LEO proxy reduces the represented two-stream
background, the exposed lunar proxy couples strongly to the prompt
$\gamma$ response, and quiet Sun--Earth L2 couples to the sparse
proton-activation response. Whether those tendencies survive in a mission
design depends on the omitted platform material, directional field, activation
history, and observing schedule. In particular, SAA exposure can leave delayed
background after a LEO instrument returns to its science orbit, and secondary
neutrons can still be produced inside an L2 payload even when no external
planetary-neutron component is specified.

\subsection{Scope and next steps}
\label{subsec:limitations}

The mission values are statistical projections for the detector--cryostat
geometries and a synthetic trajectory that remains continuously on source at
$45^\circ$ elevation. They use 50 keV offline thresholds for each geometry's
active channels (BGO plus plastic scintillator in model A and BGO in model B),
a detector-level topology test, all eight prompt families, and all eight
activation families. Background from the Laue optics and balloon platform lies
outside the detector--cryostat budget. Radiation-field, cross-section,
calibration, reconstruction, and pointing systematics are also outside the
propagated statistical uncertainty.

The model-B result is conditional on the complete configuration. Relative to
model A, its entrance geometry, local tungsten structure, BGO arrangement,
active channels, and part of the nearby material inventory change together.
The nearly equal focused effective areas establish signal comparability but do
not establish equal diffuse-photon acceptance. The present data therefore do
not support assigning the sensitivity change to the chimney displacement, BGO
layout, or any other single element in isolation.

The current statistical precision is limited mainly by the finite number and
unequal weights of selected transport events. The model-B final prompt and
broadband-$\gamma$ components contain three and twelve records, respectively,
and its cold-stage delayed estimate contains seven source records. Future
validation should increase these supports, repeat production-position sampling
to test spatial convergence, replace the synthetic trajectory with a flight
plan and visibility history, and calibrate the energy response and BGO trigger
efficiency. A spatial--spectral likelihood can then use the focused spot and
control regions while profiling background nuisance parameters
\cite{Cowan2011}. These steps would extend the present source-to-geometry logic
and quantify systematics beyond the current conditional design comparison.

\section{Conclusions}
\label{sec:conclusions}

The mass-model-A analysis identifies activated copper in the staged cold region
and nearby structures as a major source of selected delayed background. This
source map motivates the complete lateral-chimney and active-BGO configuration
of mass model B, which retains 99.75\% of the model-A selected effective area.
At day 15, the direct selected background changes from
$5.30571\times10^{-2}$ to $1.07472\times10^{-2}\cps$; model B gives
$S=2.35000\times10^{-3}\cps$ at $\ftoaref=2.4\times10^{-4}\phcms$. Along the
fixed-$45^\circ$, 81-node, 20 d trajectory, model B accumulates 4028.15 signal
and 18568.89 background counts, giving Gaussian and Asimov significances of
29.561 and 28.578. The corresponding top-of-atmosphere $3\sigma$ thresholds are
$(2.436\pm0.223)\times10^{-5}$ and
$(2.445\pm0.223)\times10^{-5}\phcms$, and the Gaussian threshold is lower by a
factor of 2.227 than for model A. This improvement is a conditional comparison
of the two complete detector--cryostat designs rather than a single-component
ablation.

The response-weighted two-stream extension gives screening indices of 0.483
for the 530\,km non-SAA LEO proxy, 3.089 for the exposed lunar proxy, and
1.570--2.223 across the two quiet Sun--Earth L2 modulation endpoints. The
transfer does not map the three-record prompt-positron residual and does not
include target-specific structures, activation histories, or common-time
closure. These values therefore set priorities for matched environmental
simulations rather than performance forecasts.

\section*{Data and software availability}
Before publication, a versioned public repository or institutional archive should provide the geometry and source definitions, stochastic-analysis conventions, analysis and figure-generation scripts, LaTeX source, and derived validation products needed to reproduce the tables and figures in this paper.

\section*{Declarations}
\textbf{Funding} To be completed before journal submission.

\textbf{Competing interests} To be completed before journal submission.

\textbf{Author contributions} To be completed before journal submission.

\textbf{Ethics approval and consent to participate} Not applicable.

\begin{thebibliography}{99}

\bibitem{Prantzos2011}
N. Prantzos, C. Boehm, A.M. Bykov, R. Diehl, K. Ferriere, N. Guessoum, P. Jean, J. Knoedlseder, A. Marcowith, I.V. Moskalenko, A. Strong, G. Weidenspointner, The 511 keV emission from positron annihilation in the Galaxy, Rev. Mod. Phys. 83 (2011) 1001--1056, doi:10.1103/RevModPhys.83.1001.

\bibitem{Jean2006}
P. Jean, J. Knoedlseder, V. Lonjou, M. Allain, J.-P. Roques, G.K. Skinner, B.J. Teegarden, G. Vedrenne, P. von Ballmoos, B. Cordier, P. Caraveo, R. Diehl, P. Durouchoux, P. Leleux, R. Schanne, G. Weidenspointner, Spectral analysis of the Galactic $e^+e^-$ annihilation emission, Astron. Astrophys. 445 (2006) 579--589, doi:10.1051/0004-6361:20053765.

\bibitem{Kierans2020}
C.A. Kierans, S.E. Boggs, A. Zoglauer, A.W. Lowell, C. Sleator, J. Beechert, T.J. Brandt, P. Jean, H. Lazar, J. Roberts, T. Siegert, J.A. Tomsick, P. von Ballmoos, Detection of the 511 keV Galactic positron annihilation line with COSI, Astrophys. J. 895 (2020) 44, doi:10.3847/1538-4357/ab89a9.

\bibitem{Siegert2016AA}
T. Siegert, R. Diehl, G. Khachatryan, M.G.H. Krause, F. Guglielmetti, J. Greiner, A.W. Strong, X. Zhang, Gamma-ray spectroscopy of positron annihilation in the Milky Way, Astron. Astrophys. 586 (2016) A84, doi:10.1051/0004-6361/201527510.

\bibitem{Siegert2016V404}
T. Siegert, R. Diehl, J. Greiner, M.G.H. Krause, A.M. Beloborodov, M. Cadolle Bel, F. Guglielmetti, J. Rodriguez, A.W. Strong, X. Zhang, Positron annihilation signatures associated with the outburst of the microquasar V404 Cygni, Nature 531 (2016) 341--343, doi:10.1038/nature16978.

\bibitem{Knodlseder2005AllSky}
J. Knoedlseder, P. Jean, V. Lonjou, G. Weidenspointner, N. Guessoum, R. Diehl, W. Gillard, M.J. Harris, G.K. Skinner, P. von Ballmoos, G. Vedrenne, P. Caraveo, B. Cordier, V. Schoenfelder, B.J. Teegarden, The all-sky distribution of 511 keV electron-positron annihilation emission, Astron. Astrophys. 441 (2005) 513--532, doi:10.1051/0004-6361:20042063.

\bibitem{Yoneda2025}
H. Yoneda, T. Siegert, S. Mittal, Imaging the positron annihilation line with 20 years of INTEGRAL/SPI observations, arXiv:2509.01066, 2025.

\bibitem{Siegert2020COSI}
T. Siegert, S.E. Boggs, J.A. Tomsick, A.C. Zoglauer, C.A. Kierans, C.C. Sleator, J. Beechert, T.J. Brandt, P. Jean, H. Lazar, A.W. Lowell, J.M. Roberts, P. von Ballmoos, Imaging the 511 keV positron annihilation sky with COSI, Astrophys. J. 897 (2020) 45, doi:10.3847/1538-4357/ab9607.

\bibitem{Siegert2019Kinematics}
T. Siegert, R. Diehl, G. Khachatryan, M.G.H. Krause, F. Guglielmetti, J. Greiner, A.W. Strong, X. Zhang, Constraints on positron annihilation kinematics in the inner Galaxy, Astron. Astrophys. 627 (2019) A126, doi:10.1051/0004-6361/201833856.

\bibitem{DeCesare2011}
G. De Cesare, Searching for the 511 keV annihilation line from galactic compact objects with the IBIS gamma ray telescope, Astron. Astrophys. 531 (2011) A56, doi:10.1051/0004-6361/201116516.

\bibitem{Barriere2009Laue}
N.M. Barriere, L. Natalucci, N. Abrosimov, P. von Ballmoos, P. Bastie, P. Courtois, M. Jentschel, J. Knodlseder, J. Rousselle, P. Ubertini, Soft gamma-ray optics: new Laue lens design and performance estimates, arXiv:0910.0488, 2009.

\bibitem{Barriere2011LauePrototype}
N.M. Barriere, J.A. Tomsick, S.E. Boggs, A. Lowell, P. von Ballmoos, Developing a second generation Laue lens prototype: high reflectivity crystals and accurate assembly, arXiv:1111.6700, 2011.

\bibitem{Weidenspointner2006MAX}
G. Weidenspointner, C.B. Wunderer, N. Barriere, A. Zoglauer, P. von Ballmoos, Monte Carlo study of detector concepts for the MAX Laue Lens gamma-ray telescope, arXiv:astro-ph/0603152, 2006.

\bibitem{Shirazi2023}
F. Shirazi, E. Gau, M.A. Hossen, D. Becker, D. Schmidt, D. Swetz, D. Bennett, D. Braun, F. Kislat, J. Gard, J. Mates, J. Weber, N.R. Cavero, S. Chun, L. Lisalda, A. West, B. Dev, F. Ferrer, R. Bose, J. Ullom, H. Krawczynski, The 511-CAM Mission: A pointed 511 keV gamma-ray telescope with a focal plane detector made of stacked transition edge sensor microcalorimeter arrays, J. Astron. Telesc. Instrum. Syst. 9 (2023) 024006, doi:10.1117/1.JATIS.9.2.024006.

\bibitem{Zhang2022TESGamma}
S. Zhang, J.-K. Xia, T. Sun, et al., Transition edge sensor based detector: from X-ray to gamma-ray, arXiv:2204.00010, 2022.

\bibitem{Hu2026}
K. Hu, M. Fritts, D. Becker, D. Schmidt, F. Zhou, A. Dasgupta, F. Kislat, M. Keller, S. Chun, A.G. Detoito, E. Gau, X. Jin, D. Bennett, D. Braun, J. Gard, J.A.B. Mates, J. Nobles, D. Radomski, N. Rodriguez Cavero, R. Snodgrass, D. Swetz, J. Ullom, S. Vengalil Menon, J. Weber, K. Wimalasena, H. Krawczynski, Design of a mission to measure the shape and substructure of the 511 keV gamma-ray line from the center of the Milky Way, J. Astron. Telesc. Instrum. Syst. 12 (2026) 025002, doi:10.1117/1.JATIS.12.2.025002.

\bibitem{Gallego2025Balloon}
S. Gallego, U. Oberlack, J. Lommler, C.M. Karwin, A. Zoglauer, P. Jean, P. von Ballmoos, C. Kierans, C.C. Sleator, J.A. Tomsick, S.E. Boggs, Bottom-up background simulations of the 2016 COSI balloon flight, Astrophys. J. 986 (2025) 116, doi:10.3847/1538-4357/add6a0.

\bibitem{Tian2026DIXE}
R. Tian, J. Mao, J. Liu, H. Jin, W. Cui, Simulation of non-X-ray background for the DIffuse X-ray Explorer (DIXE) mission, arXiv:2604.13569, 2026.

\bibitem{Caputo2017}
R. Caputo, F. Kislat, J. Racusin, on behalf of the AMEGO Team, AMEGO: Simulations of the instrument performance, Proc. 35th International Cosmic Ray Conference, PoS(ICRC2017) 783 (2018), doi:10.22323/1.301.0783.

\bibitem{Gottardi2021TESReview}
L. Gottardi, K. Nagayoshi, A review of X-ray microcalorimeters based on superconducting transition edge sensors for astrophysics and particle physics, Applied Sciences 11 (2021) 3793, doi:10.3390/app11093793.

\bibitem{Xie2026AlMn}
L. Xie, Y. Zhang, Z. Li, Z. Liu, S. Shu, J. Zhou, X. Li, H. Li, H. Gao, Y. Gu, X. Lu, Y. Zhao, C. Liu, Beyond one-thousandth energy resolution with an AlMn TES detector, arXiv:2602.11728, 2026.

\bibitem{Vavagiakis2017AlMnMagnetic}
E.M. Vavagiakis, S.W. Henderson, K. Zheng, H.-M. Cho, N.F. Cothard, B. Dober, S.M. Duff, P.A. Gallardo, G. Hilton, J. Hubmayr, K.D. Irwin, B.J. Koopman, D. Li, F. Nati, M.D. Niemack, C.D. Reintsema, S. Simon, J.R. Stevens, A. Suzuki, B. Westbrook, Magnetic sensitivity of AlMn TESes and shielding considerations for next-generation CMB surveys, arXiv:1710.08456, 2017.

\bibitem{NISTXCOM}
J.H. Hubbell, S.M. Seltzer, XCOM: Photon Cross Section Database, NIST Standard Reference Database 8 (XGAM), National Institute of Standards and Technology, doi:10.18434/T48G6X.

\bibitem{Sturner2003}
S.J. Sturner, C.R. Shrader, G. Weidenspointner, B.J. Teegarden, D. Attie, B. Cordier, R. Diehl, C. Ferguson, P. Jean, A. von Kienlin, P. Paul, F. Sanchez, S. Schanne, P. Sizun, G. Skinner, C.B. Wunderer, Monte Carlo simulations and generation of the SPI response, Astron. Astrophys. 411 (2003) L81--L84, doi:10.1051/0004-6361:20031171.

\bibitem{Agostinelli2003}
S. Agostinelli, J. Allison, K. Amako, et al., Geant4---a simulation toolkit, Nucl. Instrum. Methods Phys. Res. A 506 (2003) 250--303, doi:10.1016/S0168-9002(03)01368-8.

\bibitem{Allison2016}
J. Allison, K. Amako, J. Apostolakis, et al., Recent developments in Geant4, Nucl. Instrum. Methods Phys. Res. A 835 (2016) 186--225, doi:10.1016/j.nima.2016.06.125.

\bibitem{SanchezDelRio2011XOP}
M. Sanchez del Rio, R.J. Dejus, XOP v2.4: recent developments of the x-ray optics software toolkit, Proc. SPIE 8141 (2011) 814115, doi:10.1117/12.893911.

\bibitem{Guan2023BraggReflection}
F. Guan, M. Asai, D.A. Bartkoski, M. Kleckner, Z. Harel, M. Salehpour, Adding the X-ray Bragg reflection physical process in crystal to the Geant4 Monte Carlo simulation toolkit, part I: reflection from a crystal slab, Precision Radiation Oncology 7(1) (2023) 59--66.

\bibitem{Reiazi2025G4BraggReflection}
R. Reiazi, D. Lee, D. Bartkoski, M. Kleckner, M. Salehpour, G4BraggReflection for accurate modeling of Bragg reflection in perfect and mosaic crystals, J. Phys. D: Appl. Phys. 58 (2025) 295102, doi:10.1088/1361-6463/aded1d.

\bibitem{Zoglauer2006}
A. Zoglauer, R. Andritschke, F. Schopper, MEGAlib: The Medium Energy Gamma-ray Astronomy Library, New Astron. Rev. 50 (2006) 629--632, doi:10.1016/j.newar.2006.06.049.

\bibitem{Sato2015EXPACS}
T. Sato, Analytical model for estimating terrestrial cosmic ray fluxes nearly anytime and anywhere in the world: extension of PARMA/EXPACS, PLoS ONE 10 (2015) e0144679, doi:10.1371/journal.pone.0144679.

\bibitem{Sato2016PARMA4}
T. Sato, Analytical model for estimating the zenith angle dependence of terrestrial cosmic ray fluxes, PLoS ONE 11 (2016) e0160390, doi:10.1371/journal.pone.0160390.

\bibitem{Kondev2021NUBASE}
F.G. Kondev, M. Wang, W.J. Huang, S. Naimi, G. Audi, The NUBASE2020 evaluation of nuclear physics properties, Chinese Phys. C 45 (2021) 030001, doi:10.1088/1674-1137/abddae.

\bibitem{Kierans2017}
C.A. Kierans, S.E. Boggs, J.-L. Chiu, A. Lowell, C. Sleator, J.A. Tomsick, A. Zoglauer, M. Amman, H.-K. Chang, C.-H. Tseng, C.-Y. Yang, C.-H. Lin, P. Jean, P. von Ballmoos, The 2016 Super Pressure Balloon flight of the Compton Spectrometer and Imager, arXiv:1701.05558, 2017.

\bibitem{Karwin2023}
C.M. Karwin, T. Siegert, J. Beechert, J.A. Tomsick, T.A. Porter, M. Negro, C. Kierans, M. Ajello, I. Martinez Castellanos, A. Shih, A. Zoglauer, S.E. Boggs, Probing the Galactic diffuse continuum emission with COSI, arXiv:2310.12206, 2023.

\bibitem{Cumani2019LEO}
P. Cumani, M. Hernanz, J. Kiener, V. Tatischeff, A. Zoglauer,
Background for a gamma-ray satellite on a low-Earth orbit,
Exp. Astron. 47 (2019) 273--302,
doi:10.1007/s10686-019-09624-0.

\bibitem{Dobynde2021REDMoon}
M.I. Dobynde, J. Guo,
Radiation environment at the surface and subsurface of the Moon:
model development and validation,
J. Geophys. Res. Planets 126 (2021) e2021JE006930,
doi:10.1029/2021JE006930.

\bibitem{Lotti2021Athena}
S. Lotti et al.,
Review of the particle background of the Athena X-IFU instrument,
Astrophys. J. 909 (2021) 111,
doi:10.3847/1538-4357/abd94c.

\bibitem{Usoskin2005}
I.G. Usoskin, K. Alanko-Huotari, G.A. Kovaltsov, K. Mursula,
Heliospheric modulation of cosmic rays: monthly reconstruction for 1951--2004,
J. Geophys. Res. Space Phys. 110 (2005) A12108,
doi:10.1029/2005JA011250.

\bibitem{Kuznetsov2017}
N.V. Kuznetsov, H. Popova, M.I. Panasyuk,
Empirical model of long-time variations of galactic cosmic ray particle fluxes,
J. Geophys. Res. Space Phys. 122 (2017) 1463--1472,
doi:10.1002/2016JA022920.

\bibitem{Cowan2011}
G. Cowan, K. Cranmer, E. Gross, O. Vitells, Asymptotic formulae for likelihood-based tests of new physics, Eur. Phys. J. C 71 (2011) 1554, doi:10.1140/epjc/s10052-011-1554-0.

\end{thebibliography}
\end{document}
